%Version 3 December 2023
% See section 11 of the User Manual for version history
%
%%%%%%%%%%%%%%%%%%%%%%%%%%%%%%%%%%%%%%%%%%%%%%%%%%%%%%%%%%%%%%%%%%%%%%
%%                                                                 %%
%% Please do not use \input{...} to include other tex files.       %%
%% Submit your LaTeX manuscript as one .tex document.              %%
%%                                                                 %%
%% All additional figures and files should be attached             %%
%% separately and not embedded in the \TeX\ document itself.       %%
%%                                                                 %%
%%%%%%%%%%%%%%%%%%%%%%%%%%%%%%%%%%%%%%%%%%%%%%%%%%%%%%%%%%%%%%%%%%%%%

%%\documentclass[referee,sn-basic]{sn-jnl}% referee option is meant for double line spacing

%%=======================================================%%
%% to print line numbers in the margin use lineno option %%
%%=======================================================%%

% \documentclass[lineno,sn-basic]{sn-jnl}% Basic Springer Nature Reference Style/Chemistry Reference Style

%%======================================================%%
%% to compile with pdflatex/xelatex use pdflatex option %%
%%======================================================%%

%%\documentclass[pdflatex,sn-basic]{sn-jnl}% Basic Springer Nature Reference Style/Chemistry Reference Style


%%Note: the following reference styles support Namedate and Numbered referencing. By default the style follows the most common style. To switch between the options you can add or remove “Numbered” in the optional parenthesis. 
%%The option is available for: sn-basic.bst, sn-vancouver.bst, sn-chicago.bst%  
 
%%\documentclass[pdflatex,sn-nature]{sn-jnl}% Style for submissions to Nature Portfolio journals
%%\documentclass[pdflatex,sn-basic]{sn-jnl}% Basic Springer Nature Reference Style/Chemistry Reference Style
%\documentclass[pdflatex,sn-mathphys-num]{sn-jnl}% Math and Physical Sciences Numbered Reference Style 
%\documentclass[pdflatex,sn-mathphys-ay]{sn-jnl}% Math and Physical Sciences Author Year Reference Style
%%\documentclass[pdflatex,sn-aps]{sn-jnl}% American Physical Society (APS) Reference Style
%%\documentclass[pdflatex,sn-vancouver,Numbered]{sn-jnl}% Vancouver Reference Style
%%\documentclass[pdflatex,sn-apa]{sn-jnl}% APA Reference Style 
\documentclass[pdflatex,sn-chicago]{sn-jnl}% Chicago-based Humanities Reference Style

%%%% Standard Packages
%%<additional latex packages if required can be included here>

\usepackage{graphicx}%
\usepackage{multirow}%
\usepackage{amsmath,amssymb,amsfonts}%
\usepackage{amsthm}%
\usepackage{lmodern}  % Scalable fonts for better support across sizes
\usepackage{mathrsfs}  % Since you're already using this for mathematical symbols
\usepackage[title]{appendix}%
\usepackage{xcolor}%
\usepackage{textcomp}%
\usepackage{manyfoot}%
\usepackage{booktabs}%
\usepackage{algorithm}%
\usepackage{algorithmicx}%
\usepackage{algpseudocode}%
\usepackage{listings}%
\usepackage{lmodern}  % For Latin Modern fonts (more scalable)
\usepackage{newtxmath}  % For enhanced math font support
\usepackage{hyperref}
\usepackage{lineno}
\linenumbers
\hypersetup{
    bookmarksdepth=section  % or any desired level
}
\usepackage{orcidlink}

%%%%

%%%%%=============================================================================%%%%
%%%%  Remarks: This template is provided to aid authors with the preparation
%%%%  of original research articles intended for submission to journals published 
%%%%  by Springer Nature. The guidance has been prepared in partnership with 
%%%%  production teams to conform to Springer Nature technical requirements. 
%%%%  Editorial and presentation requirements differ among journal portfolios and 
%%%%  research disciplines. You may find sections in this template are irrelevant 
%%%%  to your work and are empowered to omit any such section if allowed by the 
%%%%  journal you intend to submit to. The submission guidelines and policies 
%%%%  of the journal take precedence. A detailed User Manual is available in the 
%%%%  template package for technical guidance.
%%%%%=============================================================================%%%%

\theoremstyle{plain}
\newtheorem{theorem}{Theorem}

\theoremstyle{definition}
\newtheorem{proposition}{Proposition}

\theoremstyle{remark}
\newtheorem{example}{Example}
\newtheorem{remark}{Remark}

\theoremstyle{definition}
\newtheorem{definition}{Definition}

\raggedbottom
%%\unnumbered% uncomment this for unnumbered level heads

\begin{document}

\title[Article Title]{Energy Patterns of Southwestern Atlantic Cyclones: Evidence for Barotropic Instability}

%%=============================================================%%
%% GivenName	-> \fnm{Joergen W.}
%% Particle	-> \spfx{van der} -> surname prefix
%% FamilyName	-> \sur{Ploeg}
%% Suffix	-> \sfx{IV}
%% \author*[1,2]{\fnm{Joergen W.} \spfx{van der} \sur{Ploeg} 
%%  \sfx{IV}}\email{iauthor@gmail.com}
%%=============================================================%%

\author*[1,2,3]{\fnm{Danilo Couto} \sur{de Souza}}\email{danilo.oceano@gmail.com\orcidlink{0000-0003-4121-7583}}

\author[1]{\fnm{Pedro Leite da} \sur{Silva Dias}}\email{pldsdias@gmail.com\orcidlink{0000-0002-4051-2962}}

\author[4]{\fnm{Carolina Barnez} \sur{Gramcianinov}}\email{cbgramcianinov@gmail.com\orcidlink{0000-0002-3919-5226}}

\author[1]{\fnm{Ricardo} \sur{Camargo}}\email{ricamarg@usp.br\orcidlink{0000-0002-9425-5391}}

\affil*[1]{\orgdiv{Institute of Astronomy, Geophysics and Atmospheric Sciences}, \orgname{São Paulo University}, \orgaddress{\street{Rua do Matão, 266}, \city{São Paulo}, \postcode{05508-090}, \state{São Paulo}, \country{Brazil}}}


\affil[2]{IRB(P\&D), Av. República do Chile, 330 - 4o andar - Centro, Rio de Janeiro - RJ, 20031-170}

\affil[3]{Brazilian Center of Risk and Resilience Studies, R. Equador, 335 - Santo Cristo, Rio de Janeiro - RJ, 20220-410}

\affil[4]{\orgdiv{Institute for Coastal Systems Analysis and Modeling}, \orgname{Helmholtz-Zentrum Hereon}, \orgaddress{\street{Max-Planck-Straße, 1}, \city{Geesthacht}, \postcode{21502}, \state{Schleswig-Holstein}, \country{Germany}}}


\abstract{Lorem ipsum}


\keywords{Cyclone energetics, Baroclinic conversion, Barotropic conversion, Latent heat release, South Atlantic cyclones}

%%\pacs[JEL Classification]{D8, H51}

%%\pacs[MSC Classification]{35A01, 65L10, 65L12, 65L20, 65L70}

\maketitle

\section*{Acknowledgements}

The first author would like to thank the Laboratory of Applied Meteorology for Regional Meteorological Systems (MASTER) for providing the data processing infrastructure. Special thanks are extended to Jean Peres and Djalma Vieira for their invaluable technical assistance and support in the laboratory, which contributed significantly to the success of this research. CBG acknowledges the Helmholtz European Partnership “Research Capacity Building for Healthy, Productive and Resilient Seas” (SEA-ReCap, grant no. PIE-0025).

\section{Introduction}\label{sec1}

Extratropical cyclones are key components of the midlatitude storm tracks and major drivers of synoptic variability, acting as transient eddies that transport heat, moisture, and momentum poleward. By organizing fronts and modulating the occurrence of precipitation and wind extremes, they exert a strong influence on regional climates and on the large-scale atmospheric energy cycle. Understanding the dynamical processes that control cyclone growth and intensification is therefore essential for interpreting climate variability and for improving numerical models used in weather prediction and climate projections.

Over South America, three cyclogenesis hotspots have been identified in the literature: eastern Argentina (ARG), the La Plata River basin (LA-PLATA), and southeastern Brazil (SE-BR) \citep{gan1991surface,hoskins2005new,gramcianinov2019properties,crespo2021potential}. Cyclones forming in these regions often undergo genesis, intensification, and maturity over or near the continent \citep{couto2024new}, substantially modulating weather across subtropical and extratropical South America and being associated with high-impact hazards such as extreme rainfall \citep[e.g.,][]{de2021ocean,de2024extreme}, intense winds \citep{cardoso2022synoptic,de2021ocean}, high sea waves \citep{da2025classification,gramcianinov2023impact}, and storm surges along coastal areas \citep{Leal2026WeatherPatterns,tecchio2024mean}. These events translate into socioeconomic losses and disruptions to local infrastructure \citep[e.g.][]{da2022climate,marengo2024impactos}, which is particularly critical given the presence of major metropolitan areas (e.g., São Paulo, Rio de Janeiro, and Buenos Aires) and strategic coastal assets such as the Port of Santos, the largest in Latin America.

The cyclone phase space (CPS) framework, which diagnoses cyclone thermal-core structure (warm vs. cold) and symmetry \citep{hart2003cyclone}, is among the most widely used objective tools for classifying cyclone types and phase transitions \citep{wood2023phase}. However, CPS is fundamentally a structural diagnostic: it tracks how the storm’s thermal and symmetry signatures evolve, but it does not explicitly diagnose the physical processes and environmental controls driving those changes (e.g., vertical shear, baroclinicity, air–sea fluxes, or PV interactions). Moreover, CPS-based classification can be ambiguous for certain hybrid cyclone types. For instance, warm seclusions can exhibit warm, relatively symmetric cores \citep[e.g.,][]{dekker2018characteristics}, yielding CPS signatures that may overlap those of tropical(-like) cyclones. In addition, some systems follow atypical CPS pathways—for example, extratropical transitions in which the core becomes cold before a marked increase in thermal asymmetry \citep{kitabatake2011climatology}. Finally, because CPS magnitudes and commonly used thresholds depend on storm size, tropospheric depth, and regional atmospheric structure, several studies have proposed adaptations of CPS parameters (e.g., analysis radius and/or pressure layers) for medicanes, Arctic cyclones/polar lows, and South Atlantic subtropical cyclones \citep{gozzo2014subtropical,comellas2021evaluation,croad2023role,saraceni2023mediterranean,miglietta2025defining}.

Some authors proposed objective cyclone classification methods based on environmental parameters \citep[e.g][]{yanase2014parameter,emanuel2025cyclops}. For instance, \citet{corner2025classification} clustered extratropical cyclones cyclone into four intensity types, that differ in terms of not only intensity but on life cycle characteristics (e.g., duration and deepening rate) and on geographical location. Meanwhile, \citet{graf2017objective} attempted to classify cyclogenesis events using over 30 environmental variables. They found that these systems represent a continuum between systems with low/high level forcing with varying importance of diabatic processes. Noteworthy, \citet{han2025system} used a machine learning approach for classifying cyclone types, although, for this method it is required previously classified systems, which can propagate biases. Here, we instead adopt a phase-resolved energetics framework, quantifying the relative contributions of baroclinic, barotropic and diabatic pathways throughout the cyclone life cycle. By focusing on energy conversions rather than solely on environmental proxies or peak intensity, our approach provides a process-based classification that distinguishes cyclones with similar intensity or background conditions but fundamentally different growth mechanisms.

Extratropical cyclone development is traditionally framed in terms of baroclinic instability, and the canonical eddy life cycle is often summarized as a phase of baroclinic growth followed by barotropic decay — i.e., a net transfer of eddy kinetic energy back to the mean flow, ultimately dissipated by friction \citep[e.g.,][]{simmons1978life,randel1985observed,ma2021lorenz,cai2007dynamical}. However, study cases using limited-area energetics and Lorenz Energy Cycle (LEC) diagnostics indicate that barotropic conversions can be positive and contribute to cyclone intensification, implying that the mean-flow kinetic energy may directly feed transient eddies during development \citep[e.g.,][]{bulic2006limited,dias2011energy,black2013universal,cavicchia2018energetics}. This pathway appears to be relatively more relevant in weakly baroclinic environments and/or in hybrid–subtropical systems, where diabatic processes and enhanced barotropic conversions may jointly promote eddy growth \citep[e.g.,][]{veiga2008analysis,pezza2014large,cavicchia2019physically}. More recently, \citet{de2025lorenz} presented a large-sample, semi-Lagrangian LEC climatology for South Atlantic cyclones, showing that barotropic conversions are ubiquitous during cyclone evolution and typically peak near maturity, and that they can even exceed baroclinic conversions for some intense clusters.

In this study, we use cyclone energetics to objectively cluster extratropical cyclones forming over the southwestern South Atlantic into distinct groups with contrasting energetic pathways. Building on the semi-Lagrangian framework of \citet{de2025lorenz}, we diagnose the Lorenz Energy Cycle (LEC) for more than 3,800 cyclones with genesis in the main cyclogenesis regions adjacent to South America. To enable a consistent comparison across storms with different lifetimes, we phase-average the energy-cycle terms over objectively defined life-cycle stages \citep{couto2024new}. The resulting multivariate energetic descriptors are then clustered to identify recurrent energy patterns. Finally, we assess the occurrence of baroclinic and barotropic instability during cyclone intensification as an independent dynamical validation of the inferred energetic regimes.

\section{Materials and methods}

\subsection{Data}\label{sec:data}

The primary data source used in this study is the European Centre for Medium-Range Weather Forecasts (ECMWF) fifth-generation reanalysis, ERA5 \citep{hersbach2020era5}. This dataset has global coverage, with a horizontal grid spacing of approximately 31 km (0.25$^\circ$). It includes 137 vertical levels, ranging from near the Earth's surface (1000 hPa) to the top of the atmosphere (1 hPa), and provides hourly data. ERA5 contains estimates of several atmospheric variables; for the energy computations, we used geopotential, zonal and meridional wind, vertical velocity on isobaric surfaces, and temperature. ERA5 reanalysis data are produced using a state-of-the-art data assimilation system that incorporates a large number of observations, including satellite and in situ measurements \citep{hersbach2020era5}.

The cyclone tracks used in this study are the same as those in \citet{gramcianinov2020analysis, couto2024new, de2025lorenz}. The tracks are derived from the "Atlantic Extratropical Cyclone Tracks Database" \citep{gramcianinov2020data}, which identifies cyclones based on central relative vorticity at 850 hPa ($\zeta_{850}$). Cyclones are tracked using the TRACK algorithm \citep{hodges1994general, hodges1995feature} and the method outlined by \citet{hoskins2002new}, spanning the period 1979--2020. Here, as in \citet{de2025lorenz}, we considered only cyclones with genesis in the ARG, LA-PLATA, and SE-BR cyclogenesis regions (Figure \ref{fig:1_tracks_genesis_frequency}a). The cyclone selection criteria and processing are the same as those employed in \citet{couto2024new, de2025lorenz}; for further details, please refer to these references. These selection criteria resulted in 6789 cyclones with genesis in these regions: 4015 in ARG (59.1\%), 1288 in LA-PLATA (18.9\%), and 1486 in SE-BR (21.8\%).

\subsection{Life Cycle Detection}
\label{sec:life_cycle}

Cyclones have variable lifetimes, being able to last for less than 24 hours to over 10 days \citep[e.g.][]{trigo2006climatology,gramcianinov2019properties}. Therefore, to normalize their duration and enable comparisons of cyclone energetics, we averaged the energy diagnostics over four distinct life cycle phases: incipient, intensification, mature, and decay, similarly to \citep{de2025lorenz}. This was performed using the open-source Python package Cyclophaser, which identifies cyclone life cycle phases from the time series of relative vorticity at the system's central position and its first derivative \citep{de2025cyclophaser}. The cyclone life cycle phases are then identified from peaks and troughs in the vorticity time series. For a detailed description of the procedure, the spatial distribution of cyclone life cycle phases in the South Atlantic, and statistics such as duration, displacement, speed, and intensity for each phase, see \citet{couto2024new}. For standardizing the cyclone life cycle, we considered only cyclones with a complete life cycle in the standard sequence — incipient, intensification, mature, and decay — resulting in 3820 systems, approximately 56\% of the original sample. Figure \ref{fig:1_tracks_genesis_frequency} summarizes the resulting tracks used in this study, along with their seasonal occurrence.


\begin{figure}[ht!]
    \centering
    \includegraphics[width=1\linewidth]{1_tracks_genesis_frequency.png}
    \caption{Cyclone tracks and genesis climatology for the three cyclogenesis regions considered in this study (1979--2020). (a) Tracks of cyclones with genesis in ARG, LA-PLATA, and SE-BR, shown in different colors. (b) Seasonal distribution of cyclone genesis by region: the outer ring indicates the fraction (and total number) of cyclones originating in each region, and the inner ring shows the seasonal breakdown (DJF, MAM, JJA, SON) within each region.}
    \label{fig:1_tracks_genesis_frequency}
\end{figure}

\subsection{Lorenz Energy Cycle Computation}

The LEC framework employed here is the same as in \citet{de2025lorenz}, which used the open-source Python application LorenzCycleToolKit \citep{de2024lorenzcycletoolkit}. A semi-Lagrangian framework was used to compute the LEC for a $15^\circ \times 15^\circ$ computational domain centered on the cyclone's central position. The energy budget equations for the zonal and eddy forms of APE and kinetic energy are expressed as:


\begin{flalign}
\frac{\partial A_Z}{\partial t} &= BA_Z - C_Z - C_A + G_Z \\
\frac{\partial A_E}{\partial t} &= BA_E - C_E + C_A + G_E \\
\frac{\partial K_Z}{\partial t} &= BK_Z + C_Z - C_K + B\Phi_Z - D_Z \\
\frac{\partial K_E}{\partial t} &= BK_E + C_E + C_K + B\Phi_E - D_E 
\end{flalign}

In these equations, APE is divided into zonal ($A_Z$) and eddy ($A_E$) components, as is kinetic energy ($K_Z$ and $K_E$, respectively). $C_A$ represents the conversion between $A_Z$ and $A_E$, $C_E$ denotes the conversion from $A_E$ to $K_E$, $C_K$ signifies the transformation from $K_E$ to $K_Z$, and $C_Z$ describes the conversion from $A_Z$ to $K_Z$. The generation of APE and the dissipation of kinetic energy are indicated by $G$ and $D$, with $G_Z$ and $G_E$ denoting the generation of $A_Z$ and $A_E$, and $D_Z$ and $D_E$ representing the dissipation of $K_Z$ and $K_E$, respectively. The terms $BA_Z$, $BA_E$, $BK_Z$, and $BK_E$ represent the fluxes of zonal and eddy APE and kinetic energy across the boundaries. The terms $B\Phi_Z$ and $B\Phi_E$ are grouped with $C_Z$ and $C_E$, respectively, because both arise from the same derivation of the $K_Z$ and $K_E$ balances. \citet{muench1965dynamics} notes the difficulty in interpreting $B\Phi_Z$ and $B\Phi_E$, indicating that they are related to the downward flux of kinetic energy (i.e., toward lower altitudes) and represent the emergence of kinetic energy at the boundaries of the computational domain. For the full mathematical formulations of these processes, we refer to \citet{muench1965dynamics} and \citet{brennan1980zonal}.

As our goal was to evaluate the energy patterns of cyclonic development in the Southwestern Atlantic, we considered only a subset of terms for analysis: $A_E$, $K_E$, $C_A$, $C_K$, $BK_E$ $BA_E$ and $G_E$, which govern eddy (cyclone) development. This selection facilitates synthesis by focusing on terms relevant to eddy development and, because we applied PCA followed by k-means clustering, reduces the number of input features, enhancing interpretability and mitigating potential issues associated with highly collinear variables.

This term selection is illustrated in Fig.~\ref{fig:2_20070643_lps_track_publication} using a diagram we refer to as the Lorenz Phase Space (LPS). As an example, we use cyclone ``20070643'' from the database, which exhibited the most negative (i.e., most intense) central relative vorticity during the 21st century. These diagrams display distinct LEC terms on the $x$- and $y$-axes; circle size scales with $K_E$, and circle color indicates $G_E$. We define two types of LPS: (1) a ``Conversion'' LPS (Fig.~\ref{fig:2_20070643_lps_track_publication}a), in which the $y$-axis shows $C_A$ and the $x$-axis shows $C_K$; and (2) an ``Import'' LPS (Fig.~\ref{fig:2_20070643_lps_track_publication}b), in which the $x$-axis represents $BK_E$ and the $y$-axis represents $BA_E$. The ``Conversion'' LPS highlights the relative importance of baroclinic and barotropic conversions, whereas the ``Import'' LPS indicates the relative importance of eddy-energy imports and exports. Plot annotations are included to aid interpretation.


\begin{figure}[ht!]
    \centering
    \includegraphics[width=1\linewidth]{2_20070643_lps_track_publication.png}
    \caption{Lorenz Phase Space (LPS) representation of the energetics of cyclone 20070643. (a) ``Conversion'' phase space, showing the barotropic conversion $C_K$ on the $x$-axis and the baroclinic conversion $C_A$ on the $y$-axis. (b) ``Import'' phase space, showing the boundary fluxes of eddy APE ($BA_E$, $x$-axis) and eddy kinetic energy ($BK_E$, $y$-axis). In (a,b), circles denote 6-hourly samples connected in time; circle size is proportional to eddy kinetic energy ($K_E$) and colors indicate eddy APE generation ($G_E$; W\,m$^{-2}$). Gray bands mark the zero lines, and annotations indicate the physical interpretation of each quadrant; letters A and Z mark the beginning and end of the trajectory. The black-outlined circle denotes the time of minimum (most negative) central relative vorticity. (c) Cyclone track, with 6-hourly positions colored by central relative vorticity at 850\,hPa ($-1 \times 10^{-5}$\,s$^{-1}$) and marker size scaled by $K_E  \times 10^{-5} J m^{-2}$; genesis and lysis are indicated by the green circle and red cross, respectively.}
    \label{fig:2_20070643_lps_track_publication}
\end{figure}

The example system had genesis at 2007-07-24 14:00 UTC in the LA-PLATA region and lysis at 2007-07-29 03:00 UTC in the Southeastern Atlantic (Fig.~\ref{fig:2_20070643_lps_track_publication}c). The LPS diagrams (based on 6-hourly data) reveal that, from a conversion perspective (Fig.~\ref{fig:2_20070643_lps_track_publication}a), the system is initially driven by both barotropic and baroclinic processes. During early intensification, it increasingly feeds the zonal circulation via barotropic conversion ($K_E \rightarrow K_Z$). From late intensification through maturity and early decay, barotropic conversions dominate the energy cycle, and during late decay both $C_A$ and $C_K$ indicate energy transfer from the eddy to the zonal circulation. From the boundary-flux perspective (Fig.~\ref{fig:2_20070643_lps_track_publication}b), the system exhibits nearly neutral boundary fluxes of $A_E$ at genesis, while $K_E$ is imported. During early intensification, there is a net influx of both forms of eddy energy. Over the subsequent evolution, $BK_E$ tends to decrease and even indicates export prior to maturity. After maturity, although there are periods with export of both forms of eddy energy, the boundary fluxes are, overall, closer to neutral. There is also a large diabatic contribution ($G_E$) during early intensification, which decreases over time. Lastly, $K_E$ increases gradually from genesis to maturity and decreases rapidly afterward.


We also investigated the contribution of each subterm associated with $C_K$, as these provide a more detailed perspective on the dynamical processes governing barotropic conversions:


\begin{equation}
\begin{split}
C_K = \int_{p_t}^{p_b} \frac{1}{g}  & \Bigg( \underbrace{ \left[ \frac{\cos\phi}{a} (u)_{\lambda} (v)_{\lambda} \frac{\partial}{\partial\phi} \left(\frac{[u]_{\lambda}}{\cos\phi}\right) \right]_{\lambda\phi}}_{\text{(A)}} + \underbrace{\left[ \frac{(v)_{\lambda}^2}{a} \frac{\partial [v]_{\lambda}}{\partial\phi}\right]_{\lambda\phi}}_{\text{(B)}} + \underbrace{\left[ \frac{\tan\phi}{a} (u)_{\lambda}^2 [v]_{\lambda}\right]_{\lambda\phi}}_{\text{(C)}} \\
& + \underbrace{ \left[(\omega)_{\lambda}  (u)_{\lambda} \frac{\partial [u]_{\lambda}}{\partial p}\right]_{\lambda\phi}}_{\text{(D)}} + \underbrace{ \left[(\omega)_{\lambda}  (v)_{\lambda} \frac{\partial [v]_{\lambda}}{\partial p}\right]_{\lambda\phi}}_{\text{(E)}} \Bigg)  dp
\end{split}
\label{eq:CK}
\end{equation}

Where \(p_t\) and \(p_b\) are the top and bottom pressure levels used for the analysis (100 hPa and 1000 hPa, respectively), \(g\) is the gravitational acceleration, \(\phi\) is longitude, \(\lambda\) is latitude, \(a\) is the Earth’s radius, \(u\) and \(v\) are the zonal and meridional wind components, and \(\omega\) is the vertical velocity in isobaric coordinates (pressure, \(p\)). The \([X]_{\lambda}\) denotes the zonal mean, while \((X)\) represents eddy-related terms. The complete mathematical expressions for all LEC terms, including the eddy–mean decomposition, treatment of lateral fluxes, pressure–work terms, and an extended methodological discussion, are provided in \citet{de2025lorenz}.

The $C_K$ subterms can be interpreted as follows. Term (A) expresses the eddy momentum flux linked to the meridional (north--south) gradient of the zonal wind. Term (B) represents the meridional flux of kinetic energy associated with the meridional wind and its meridional gradient. Term (C) describes the meridional flux of zonal kinetic energy. Term (D) captures the zonal and vertical flux associated with the vertical shear of the zonal wind, whereas term (E) accounts for the meridional and vertical flux associated with the vertical shear of the meridional wind. Among these contributions, term (A) is the most directly connected to barotropic instability, since shear instabilities are fundamentally driven by horizontal wind shear \citep{rayleigh1895stability,kuo1949dynamic,chen1985note,mak1989local,yokota2012tropical,read2020baroclinic}. The remaining terms represent secondary contributions and/or modulate the primary instability, reflecting interactions between horizontal and vertical fluxes of kinetic energy.

\subsection{Clustering Analysis}
\label{sec:clustering}

Here, we perform a cluster analysis for identifying the Energy Patterns (EP) related to cyclones in the Southwestern Atlantic. For this we used the main LEC terms related to eddy development ($A_E$, $K_E$, $C_A$, $C_K$, $BK_E$ $BA_E$ and $G_E$), as previously discussed. This analysis consisted of the three steps, which are described bellow: 1) principal component analysis (PCA) for reducing dimensionality 2) determining the best number of clusters using objective indices, 3) applying the clustering analysis.

Given that more than 3{,}800 cyclones were analysed in this study, we applied PCA to reduce the dimensionality of the LEC feature space prior to clustering. PCA constructs a set of mutually orthogonal principal components (PCs) as linear combinations of the original variables, ordered such that PC1 explains the largest fraction of the total variance, followed by PC2, and so on. Dimensionality reduction is achieved by retaining only the first $K$ PCs that explain most of the variability in the data. This step is widely adopted in meteorological pre-clustering workflows \citep[e.g.,][]{camus2014,de2022hybrid,borato2024evaluation,pozo2025characterising,zhong2025advanced}. The PCA was performed with the open-source Python library PyEOF \citep{zheng_2021}.

Let $N$ be the number of cyclones and $p=28$ the number of features. We organised the predictors into a matrix
\begin{equation}
\mathbf{X} \in \mathbb{R}^{N \times p},
\end{equation}
whose columns are the LEC terms evaluated as phase-mean values over four lifecycle phases (incipient, intensification, mature, and decay). Specifically, for each term $n=1,\ldots,7$ and phase $\phi \in \{\mathrm{ic},\mathrm{it},\mathrm{ma},\mathrm{de}\}$, a feature $T_{n,\phi}$ is included, so that
\begin{equation}
\mathbf{X} = \big[T_{1,\mathrm{ic}}, \ldots, T_{7,\mathrm{ic}},\;
               T_{1,\mathrm{it}}, \ldots, T_{7,\mathrm{it}},\;
               T_{1,\mathrm{ma}}, \ldots, T_{7,\mathrm{ma}},\;
               T_{1,\mathrm{de}}, \ldots, T_{7,\mathrm{de}}\big],
\end{equation}
yielding $7 \times 4 = 28$ features.

Prior to PCA, each feature (column) was standardised to have zero mean and unit variance using the \texttt{StandardScaler} implementation in scikit-learn \citep{pedregosa2011scikit}. Denoting by $x_{ij}$ the value for cyclone $i$ and feature $j$, we computed
\begin{align}
\mu_j &= \frac{1}{N}\sum_{i=1}^{N} x_{ij},\\
s_j  &= \sqrt{\frac{1}{N-1}\sum_{i=1}^{N} (x_{ij}-\mu_j)^2},\\
z_{ij} &= \frac{x_{ij}-\mu_j}{s_j},
\end{align}
forming the standardised matrix $\mathbf{Z}=[z_{ij}] \in \mathbb{R}^{N \times p}$ on which PCA was applied. We retained the first $K=15$ PCs, which explain approximately $90\%$ of the cumulative variance (Figure \ref{fig:S1_pca_clustering_validation}a).

After dimensionality reduction with PCA, we clustered the cyclones using the $k$-means algorithm \citep{macqueen1967}. Owing to its simplicity, versatility, and low computational cost, $k$-means is one of the most widely used clustering methods in the geosciences, although other powerful alternatives exist \citep[e.g.,][]{fachieri2025introduccao}. Let $\mathbf{Y}\in\mathbb{R}^{N\times K}$ denote the matrix of PCA scores for the $N$ cyclones retaining the first $K$ principal components, and let $\mathbf{y}_i\in\mathbb{R}^{K}$ be the $i$-th row of $\mathbf{Y}$. The $k$-means method partitions the set of points $\{\mathbf{y}_i\}_{i=1}^{N}$ into $C$ clusters by minimizing the within-cluster sum of squares (WCSS),

\begin{equation}
\mathcal{J} = \sum_{c=1}^{C}\sum_{i \in \mathcal{I}_c}\left\lVert \mathbf{y}_i - \boldsymbol{\mu}_c \right\rVert^2,
\end{equation}

where $\mathcal{I}_c$ is the index set of points assigned to cluster $c$ and $\boldsymbol{\mu}_c\in\mathbb{R}^{K}$ is the cluster centroid (i.e., the mean of the points in cluster $c$). The optimization is performed iteratively using Lloyd's algorithm: (i) assign each point to the nearest centroid in Euclidean distance, and (ii) update each centroid as the mean of the points currently assigned to it, repeating until convergence \citep{lloyd1982least}. Here, we used the scikit-learn implementation of $k$-means with $k$-means++ initialization \citep{arthur2006k}, which selects initial centroids in a distance-aware manner to improve convergence and reduce the chance of poor local minima.

Because $k$-means requires the number of clusters to be specified \emph{a priori}, we adopted an objective, multi-metric procedure to select this value. More than 100 cluster validity indices (CVIs) have been proposed in the literature \citep{arbelaitz2013extensive,hassan2024z}; here we combine five complementary CVIs that quantify cluster compactness and separation from different perspectives (Table~\ref{tab:cvi_summary}). For each candidate number of clusters $C$ (tested over a prescribed range), we computed all five indices from the $k$-means partition in the retained PCA space.

To place all CVIs on a common scale, each index $I_m(C)$ (with $m=1,\ldots,5$) was rescaled to $[0,1]$ using min--max normalization,
\begin{equation}
\widetilde{I}_m(C) = \frac{I_m(C) - \min\limits_{C} I_m(C)}{\max\limits_{C} I_m(C) - \min\limits_{C} I_m(C)}.
\end{equation}
For indices where smaller values indicate better clustering (Davies--Bouldin), we inverted the normalized score as
\begin{equation}
\widetilde{I}_{\mathrm{DB}}^{\,*}(C) = 1 - \widetilde{I}_{\mathrm{DB}}(C).
\end{equation}
Finally, we defined a composite (unweighted) agreement score as the mean of the normalized CVIs,
\begin{equation}
\overline{I}(C) = \frac{1}{5}\sum_{m=1}^{5}\widetilde{I}_m^{\,*}(C),
\end{equation}
and selected the optimal number of clusters as the maximizer of $\overline{I}(C)$. This procedure consistently indicated $C=3$ as the optimal partition for the full dataset, supporting three distinct energy patterns while reducing reliance on any single validity criterion  (Figure \ref{fig:S1_pca_clustering_validation}b). 


\begin{table}[t]
\centering
\caption{Cluster validity indices (CVIs) used to select the number of clusters $C$: Silhouette (SIL; \citealp{rousseeuw1987silhouettes}), Davies--Bouldin (DB; \citealp{davies2009cluster}), Calinski--Harabasz (CH; \citealp{calinski1974dendrite}), Score Function (SF; \citealp{saitta2008comprehensive}), and Gap statistic (GAP; \citealp{tibshirani2001estimating}). Optimization indicates whether the index should be maximized (Max) or minimized (Min).}
\label{tab:cvi_summary}
\begin{tabular}{p{0.8cm}p{1cm}p{0.8cm}p{7.2cm}}
\toprule
\textbf{Index} & \textbf{Range} & \textbf{Opt.} & \textbf{What it measures (summary)} \\
\midrule
SIL & $[-1,\,1]$ & Max &
Balance between within-cluster cohesion and separation to the nearest neighbouring cluster; values near 1 indicate well-separated, compact clusters. \\
DB  & $[0,\,\infty)$ & Min &
Average (worst-case) similarity between each cluster and its most similar other cluster, based on within-cluster scatter and between-centroid distance; lower implies better separation/compactness. \\
CH  & $[0,\,\infty)$ & Max &
Variance ratio criterion: ratio of between-cluster dispersion to within-cluster dispersion (adjusted by sample size); higher implies denser and more separated clusters. \\
SF  & $(0,\,1)$ & Max &
Normalized measure combining compactness and separation via a bounded (sigmoid) transformation of the log-ratio between- and within-cluster distances. \\
GAP & $[0,\,\infty)$ & Max &
Difference between observed within-cluster dispersion and that expected under a null reference distribution; larger gaps indicate stronger clustering structure than random. \\
\bottomrule
\end{tabular}
\end{table}



\subsection{Diagnositc fields}
\label{sec:diagnostic_fields}

To investigate the key physical processes driving cyclonic development, we performed a composite analysis centered on cyclone centers. The analysis is restricted to the intensification phase as a proxy for the mechanisms governing cyclone development. For each cyclone, the composite was constructed over a $30^{\circ} \times 30^{\circ}$ domain centered on the cyclone center at the central time step of the intensification phase. This domain allows us to examine both the dynamics within the LEC domain ($30^{\circ} \times 30^{\circ}$) and those in the surrounding environment, thereby enabling a distinction between the dynamical processes occurring inside and outside the domain. For cyclones whose intensification phase contained an odd number of time steps, the three central time steps were used, whereas for those with an even number of time steps, the two central time steps were used.

For the composite analysis, we also excluded systems whose intensification phase lasted less than 24 hours. This was done to exclude short-lived systems and to avoid contaminating the analysis with time steps too close to the incipient or mature stages, during which the energetics exhibit distinct behavior \citep{de2025lorenz} and, therefore, likely also the dynamical fields. After this filtering, the number of analyzed cyclones decreased from 3820 to 2733. For each energy pattern, the reduction was as follows: EP1, 332 cyclones (from 444); EP2, 776 cyclones (from 979); and EP3, 1625 cyclones (from 2397).


Potential vorticity (PV) is a useful diagnostic of cyclonic development because, for adiabatic and frictionless flow, it is approximately materially conserved and combines the effects of absolute vorticity and stratification in a single quantity \citep{hoskins1985use}. In this study, PV is computed using the hydrostatic Ertel formulation in pressure coordinates \citep{bluestein1993observations}:

\begin{equation}
\mathrm{PV}
=
-g\left(
\frac{\partial u}{\partial p}\frac{\partial \theta}{\partial y}
-
\frac{\partial v}{\partial p}\frac{\partial \theta}{\partial x}
+
(\zeta + f)\frac{\partial \theta}{\partial p}
\right),
\end{equation}

where \(u\) and \(v\) are the zonal and meridional wind components, respectively, \(\theta\) is the potential temperature, \(p\) is pressure, \(g\) is the gravitational acceleration, and \(f\) is the Coriolis parameter. The relative vorticity is defined as

\begin{equation}
\zeta = \frac{\partial v}{\partial x} - \frac{\partial u}{\partial y}.
\end{equation}

All partial derivatives were evaluated using centered finite differences at interior grid points, while forward and backward one-sided finite differences were applied at the domain boundaries. The same numerical scheme was adopted for the computation of temperature advection (AdvT).

As a proxy for baroclinic instability, we used the maximum Eady growth rate (EGR), which estimates the growth rate of the most unstable baroclinic mode in the Eady model \citep{eady1949long,lindzen1980simple}. In its continuous form, EGR is given by

\begin{equation}
\mathrm{EGR} = 0.31 \frac{|f|}{N} \left| \frac{\partial \mathbf{V}}{\partial z} \right|,
\end{equation}

where \(f\) is the Coriolis parameter, \(N\) is the Brunt--V\"ais\"al\"a frequency, \(z\) is height, and \(\mathbf{V}=(u,v)\) is the horizontal wind vector. The Brunt--V\"ais\"al\"a frequency is defined as

\begin{equation}
N = \left( \frac{g}{\theta} \frac{\partial \theta}{\partial z} \right)^{1/2},
\end{equation}

where \(g\) is the gravitational acceleration and \(\theta\) is potential temperature.

As a compact diagnostic of the horizontal redistribution of upper-level kinetic energy around jet streaks and troughs, upper-level kinetic energy advection at 250 hPa was computed as

\begin{equation}
\mathrm{KE}_{\mathrm{adv}}
=
-\mathbf{V}\cdot\nabla\left(\frac{1}{2}\lvert \mathbf{V}\rvert^2\right),
\end{equation}

where \(\mathbf{V}=(u,v)\) is the horizontal wind vector at 250 hPa. Positive values of \(\mathrm{KE}_{\mathrm{adv}}\) indicate local import of kinetic energy by horizontal advection, whereas negative values indicate local export.

To assess conditions favorable to barotropic instability, we used a local Rayleigh--Kuo-type diagnostic at 250 hPa \citep{rayleigh1895stability,kuo1949dynamic}:

\begin{equation}
\mathrm{RK}
=
\beta - \frac{\partial^2 u}{\partial y^2},
\end{equation}

where \(\beta = \partial f/\partial y\), \(u\) is the zonal wind component, and \(y\) is the meridional coordinate. In classical barotropic instability theory, for a parallel basic flow \(U(y)\), a necessary condition for instability is that \(\beta - \partial^2 U/\partial y^2\) changes sign somewhere within the domain. Here, \(\mathrm{RK}\) was evaluated locally from the 250 hPa zonal wind field as a proxy for regions where the meridional gradient of absolute vorticity may reverse sign, indicating an environment favorable to barotropic instability.


Ageostrophic geopotential flux convergence (AFC) was computed at 250 hPa as a diagnostic of downstream development and local energization of the disturbance \citep[e.g.][]{orlanski1991life,orlanski1993ageostrophic}:

\begin{equation}
\mathrm{AFC}
=
-\nabla \cdot \left( \mathbf{v}_{ag}' \, \Phi' \right),
\end{equation}

where \(\mathbf{v}_{ag}'\) is the anomalous ageostrophic wind and \(\Phi'\) is the geopotential anomaly. The anomalous ageostrophic wind was defined as

\begin{equation}
\mathbf{v}_{ag}' = \mathbf{v}' - \mathbf{v}_g',
\end{equation}

where \(\mathbf{v}'\) is the horizontal wind anomaly and \(\mathbf{v}_g'\) is the geostrophic wind anomaly derived from the geopotential anomaly. Positive AFC values indicate local convergence of anomalous ageostrophic geopotential fluxes, which is consistent with downstream development and the local accumulation of
eddy kinetic energy \citep[e.g.][]{orlanski1995stages,ndarana2021downstream,pinheiro2022contributions}.

In the original downstream-development framework, perturbations are commonly defined relative to a background mean flow or to a low-frequency component of the circulation \citep[e.g.][]{pinheiro2022contributions,pinheiro2024deepening}. Here, as an approximation suitable for the composite analysis, anomalies were defined as deviations from the 30-year monthly climatological mean for the corresponding calendar month. For events spanning two different months, the background state was obtained from a weighted mean of the two monthly climatologies in order to avoid artificial discontinuities.

\subsection{Statistical Analysis}
\label{sec:stats}

Beyond the clustering procedure itself (Sect.~\ref{sec:clustering}), three further scientific questions required a dedicated statistical treatment. First, having partitioned cyclones into EP1, EP2, and EP3 on the basis of a distance-minimizing algorithm, we needed to establish whether these groups are also statistically distinguishable in terms of their LEC diagnostics and dynamical structure, rather than being an arbitrary partition of a continuum. Second, we needed to quantify how strongly cyclone energetics covary with the dynamical fields described in Sect.~\ref{sec:diagnostic_fields}, in order to connect the phase-resolved energy pathways to physically interpretable structures such as PV anomalies, baroclinicity, and upper-level flux convergence. Third, because each EP shows different occurrence frequencies through the 42-year record (Fig.~\ref{fig:6_ep_intensity_seasonality_trends}c), we needed an objective test of whether these interannual fluctuations reflect genuine long-term trends or are consistent with sampling variability alone. These three questions define the three statistical tracks described below: (i) inter-EP comparison of scalar variables, (ii) association analysis between LEC terms and dynamical descriptors, and (iii) trend analysis of EP occurrence.

The statistical analysis was performed for all 24 LEC terms and for scalar features derived from the storm-centered dynamical fields used in the composite analysis. Rather than adopting a pointwise approach, each two-dimensional field was summarized into 13 scalar descriptors computed over the inner storm-centered domain: the spatial mean, the mean absolute value, the value at the cyclone center, cardinal-sector means (north, south, east, and west), border means (north, south, east, and west), and north--south and east--west contrasts. This strategy was adopted to reduce the dimensionality of the composite fields while retaining physically interpretable measures of field intensity and structural asymmetry. In particular, the mean absolute value was included because the spatial mean of signed fields, such as temperature advection and ageostrophic flux convergence, may approach zero owing to sign cancellation between opposing sectors even when the field is locally intense. The use of cardinal sectors and zonal/meridional contrasts is also physically consistent with the LEC framework, which is itself built from zonal means and deviations from the domain mean. Combining the 24 LEC terms with the composite-derived descriptors yielded 154 scalar variables subjected to the inter-EP testing procedure described below.

\subsubsection{Testing whether the energy patterns differ: a sequential decision procedure}

For each of the 154 variables, we followed a sequential decision procedure to test whether EP1, EP2, and EP3 differ, choosing the appropriate global test and post-hoc comparison based on the distributional properties of that specific variable rather than assuming a single test applies uniformly to all of them.

The first step evaluates whether the values within each EP group are approximately Gaussian, using the Shapiro--Wilk test \citep{shapiro1965analysis}. This test compares the ordered sample values against the values expected from a normal distribution with the same mean and variance; large deviations between the two yield a small \emph{p}-value, indicating that the group departs from normality. The second step evaluates whether the three groups share a common variance, using the Brown--Forsythe modification of Levene's test \citep{levene1960robust,brown1974robust}, in which the test statistic is a one-way ANOVA \(F\) computed not on the raw values but on their absolute deviations from the group \emph{median} (rather than the mean, as in the original Levene formulation). Centering on the median makes the test more robust when the underlying distributions are skewed or heavy-tailed, as is typical of LEC terms, which are often bounded near zero for weak systems and long-tailed for intense conversions.

The outcome of these two checks determines the global test used to compare the three EPs, and, in turn, the correct post-hoc procedure and effect-size measure:
\begin{itemize}
\item \textbf{Normal in every group, homogeneous variances:} one-way ANOVA, followed if significant by Tukey's HSD pairwise test, with \(\omega^2\) as the effect size.
\item \textbf{Normal in every group, heterogeneous variances:} Welch's ANOVA (which does not assume equal variances), followed if significant by pairwise Welch \(t\)-tests with Holm's correction (used in place of the Games--Howell test, to which it is methodologically equivalent when combined with a familywise-error correction), with \(\omega^2\) as the effect size.
\item \textbf{Non-normal in at least one group:} the Kruskal--Wallis rank-based test \citep{kruskal1952use}, followed if significant by Dunn's post-hoc procedure \citep{dunn1964multiple} with Holm's correction, with \(\varepsilon^2\) as the effect size.
\end{itemize}
In the present dataset, every one of the 154 tested variables departed from normality in at least one EP group, so all inter-EP comparisons in this study follow the third (non-parametric) branch: Kruskal--Wallis as the global test and Dunn's procedure as the post-hoc test. This is not an assumption imposed \emph{a priori}: because sample sizes per EP are large (\(n \sim 330\)--2400), the Shapiro--Wilk test is also highly sensitive to small departures from normality that carry little practical consequence for symmetric, well-behaved variables; nonetheless, we retained the conservative, non-parametric route for every variable rather than judging normality case by case, since it does not assume equal variances or symmetric distributions and remains valid regardless of the (often skewed) shape of LEC terms.

\paragraph{Global test: is there a difference at all?} The Kruskal--Wallis statistic tests the null hypothesis that all EPs are draws from the same underlying distribution against the alternative that at least one EP is stochastically larger or smaller than the others. All \(N\) observations from the three EPs are pooled and jointly ranked (tied values receive the average of their ranks), and the statistic compares the mean rank of each group to the mean rank expected under the null hypothesis:
\begin{equation}
H = \frac{12}{N(N+1)} \sum_{c=1}^{3} \frac{R_c^2}{n_c} - 3(N+1),
\label{eq:kw}
\end{equation}
where \(n_c\) is the number of cyclones in EP \(c\), \(R_c\) is the sum of ranks in EP \(c\), and \(N = \sum_c n_c\). Under the null hypothesis, \(H\) approximately follows a \(\chi^2\) distribution with 2 degrees of freedom (for 3 groups); large \(H\) indicates that at least one EP occupies systematically higher or lower ranks than the others. In other words, if EP1's values for a given LEC term tend to fall consistently above those of EP2 and EP3 once all cyclones are sorted together, \(H\) will be large regardless of whether the underlying distributions are Gaussian, skewed, or multimodal — this is precisely the robustness that motivates using a rank-based test on LEC terms. A significant Kruskal--Wallis result only establishes that the three EPs are \emph{not} all identical; it does not identify which pair (or pairs) of EPs differ, which requires the post-hoc step described next.

\paragraph{Post-hoc test: which pairs differ?} When the global test was significant, we performed all three pairwise contrasts (EP1 vs.\ EP2, EP1 vs.\ EP3, EP2 vs.\ EP3) using Dunn's test \citep{dunn1964multiple}, which compares the mean rank of each pair of groups using the same joint ranking employed in the Kruskal--Wallis test:
\begin{equation}
z_{ij} = \frac{\bar{R}_i - \bar{R}_j}{\sqrt{\left[\dfrac{N(N+1)}{12} - \dfrac{\sum_t (t^3-t)}{12(N-1)}\right]\left(\dfrac{1}{n_i}+\dfrac{1}{n_j}\right)}},
\label{eq:dunn}
\end{equation}
where \(\bar{R}_i\) and \(\bar{R}_j\) are the mean ranks of EPs \(i\) and \(j\), and the summation over \(t\) is a tie-correction term. The resulting \(z_{ij}\) is converted to a two-sided \emph{p}-value from the standard normal distribution and then adjusted for multiple comparisons within that variable's three contrasts using Holm's sequential procedure \citep{holm1979simple} (Sect.~\ref{sec:stats-mcp}). The distinction between the global and post-hoc steps is central to how results are reported throughout this study: a significant Kruskal--Wallis test establishes only that ``the three groups are not all equal''; it takes the post-hoc Dunn contrasts to establish, for example, that ``EP1 has significantly larger barotropic conversion (\(C_K\)) than both EP2 and EP3, while EP2 and EP3 are not significantly different from each other.''

\subsubsection{Effect size: how large is the difference, not just whether it exists}
\label{sec:stats-effect}

A \emph{p}-value indicates whether a difference between EPs is detectable given the data, but it says nothing about its magnitude. This distinction matters greatly here because EP sample sizes range from a few hundred to over two thousand cyclones (Sect.~\ref{sec:composites}), and with such sample sizes even physically trivial differences in central tendency can produce very small \emph{p}-values. We therefore report an effect-size measure alongside every global and pairwise comparison, and use these, together with physical plausibility, as the primary basis for interpretation.

For the global Kruskal--Wallis comparisons, we used \(\varepsilon^2\), an analogue of the proportion-of-variance-explained statistics used in parametric ANOVA, computed as
\begin{equation}
\varepsilon^2 = \frac{H - k + 1}{N - k},
\label{eq:epsilon2}
\end{equation}
where \(H\) is the Kruskal--Wallis statistic (Eq.~\ref{eq:kw}), \(k=3\) is the number of EPs, and \(N\) is the total sample size. \(\varepsilon^2\) ranges from 0 (no separation between EP rank distributions) to 1 (complete separation), and we follow the conventional bins of \citet{rea1992designing} to describe its magnitude: negligible (\(<0.01\)), small (0.01--0.06), medium (0.06--0.14), and large (\(\ge 0.14\)).

For the pairwise Dunn contrasts, we used the rank-biserial correlation coefficient,
\begin{equation}
r_{rb} = 1 - \frac{2U}{n_i\,n_j},
\label{eq:rankbiserial}
\end{equation}
where \(U\) is the Mann--Whitney \(U\) statistic \citep{mann1947test} computed between the two EPs being compared. \(r_{rb}\) ranges from \(-1\) to \(+1\): a value of \(+1\) means that every value in group \(i\) exceeds every value in group \(j\) (complete separation, group \(i\) larger); \(-1\) indicates the reverse; and 0 indicates that the two groups are statistically indistinguishable in rank. In other words, \(r_{rb}\) is directly interpretable as the difference between the probability that a randomly drawn cyclone from EP \(i\) exceeds a randomly drawn cyclone from EP \(j\), and the reverse probability — a genuinely probabilistic effect size, not merely a rescaled test statistic. Following standard conventions \citep{cohen1988statistical}, we treat \(|r_{rb}| < 0.10\) as negligible, \(0.10\text{--}0.30\) as small, \(0.30\text{--}0.50\) as medium, and \(\ge 0.50\) as large; these are the bins used in the discrete color scales of Figs.~\ref{fig:S2_pairwise_effectsize_lec_terms} and \ref{fig:S4_pairwise_effectsize_composite_scalars}. The sign of \(r_{rb}\) (equivalently, of the median difference) indicates which EP has the larger value; for example, a positive EP1-vs-EP2 contrast for \(C_K\) means that EP1 systems have systematically larger (more positive) barotropic conversion than EP2 systems.

\subsubsection{Statistical significance, sample size, and physical relevance}
\label{sec:stats-pvalue}

Because this study combines a large cyclone sample with many tested variables, we adopt an explicit interpretive protocol that keeps statistical significance and physical importance conceptually separate throughout the Results.

A \emph{p}-value is the probability, computed under the null hypothesis of no difference (or no association), of observing a test statistic at least as extreme as the one obtained from the data. When \(p < \alpha\) (here \(\alpha = 0.05\) throughout), the observed pattern would be unlikely if the null hypothesis were exactly true, and we treat this as evidence against it. Critically, \(p<\alpha\) does \emph{not} mean that the alternative hypothesis is true, that the effect is large, or that it is physically important — it only means that pure chance is an unlikely explanation for the observed pattern, given the null model and the amount of data available. Conversely, \(p \ge \alpha\) does not prove the null hypothesis; it only means the data do not provide sufficiently strong evidence against it. With hundreds to thousands of cyclones per group, statistical tests in this study are highly powered: even small, physically marginal differences in central tendency (for the Kruskal--Wallis/Dunn comparisons) or weak monotonic covariability (for the correlation analysis, Sect.~\ref{sec:stats-corr}) can register as ``significant'' at conventional thresholds. For this reason, adjusted \emph{p}-values are used in this study primarily as a screening step — to flag which variables merit closer examination — while effect sizes (\(\varepsilon^2\), \(r_{rb}\), and correlation magnitudes) and consistency with the composite dynamical fields (Sect.~\ref{sec:composites}) are used as the primary basis for physical interpretation. A result is treated as evidentially strong in this study only when statistical significance, a non-negligible effect size, and physical plausibility agree; results for which only significance is reached (typical of variables with small effect size but very large \(N\)) are discussed, when at all, as exploratory rather than as robust findings.

\subsubsection{Controlling for multiple comparisons}
\label{sec:stats-mcp}

Testing 154 variables, each with up to three pairwise contrasts, substantially inflates the chance of at least one false-positive result if left uncorrected: even under a scenario in which none of the true underlying differences exist, roughly 5\% of 154 independent tests — about eight — would be expected to appear ``significant'' by chance alone at \(\alpha=0.05\). We therefore apply multiple-comparison correction at two distinct levels, which serve different purposes.

Within each variable, the three pairwise Dunn contrasts (EP1--EP2, EP1--EP3, EP2--EP3) are corrected using Holm's sequential procedure \citep{holm1979simple}: the three raw \emph{p}-values are sorted from smallest to largest, and each is compared against a progressively less strict threshold (\(\alpha/3\) for the smallest, \(\alpha/2\) for the next, and \(\alpha\) for the largest), stopping and accepting all remaining contrasts as non-significant as soon as one comparison fails. This controls the familywise error rate — the probability of making \emph{any} false-positive claim within that variable's three contrasts.

Across variables, the raw \emph{p}-values from the global Kruskal--Wallis tests are further adjusted \emph{within each analysis block} — the block of 24 LEC terms and the block of composite-derived dynamical descriptors are corrected separately, since they represent conceptually distinct families of hypotheses — using the Benjamini--Hochberg false discovery rate (FDR) procedure \citep{benjamini1995controlling}. Given \(m\) ordered \emph{p}-values \(p_{(1)} \le p_{(2)} \le \dots \le p_{(m)}\), the procedure finds the largest rank \(i^{*}\) such that
\begin{equation}
p_{(i^{*})} \le \frac{i^{*}}{m}\,q,
\label{eq:bh}
\end{equation}
and declares all variables with rank \(\le i^{*}\) significant at FDR level \(q\) (here \(q=0.05\)). Unlike Holm's procedure, which guards against declaring \emph{any} false positive, the Benjamini--Hochberg procedure controls the \emph{expected proportion} of false positives among all variables declared significant — an appropriate criterion when screening a large number of variables for the strongest, physically interpretable candidates, rather than testing a small number of pre-specified contrasts. In practice, this means that after FDR correction, an adjusted \emph{p}-value (sometimes called a \(q\)-value) below 0.05 for a given LEC term or dynamical descriptor implies that, among all variables flagged as significant at that threshold, no more than roughly 5\% are expected to be false discoveries.

\subsubsection{Association between cyclone energetics and dynamical structure}
\label{sec:stats-corr}

To assess how LEC terms covary with the dynamical fields described in Sect.~\ref{sec:composites} (Sect.~\ref{sec:correlation}), we computed, for each pair of a canonical LEC term and a scalar descriptor derived from a dynamical field, both Pearson's product-moment correlation coefficient \citep{pearson1920notes},
\begin{equation}
r = \frac{\sum_{i=1}^{n}(x_i-\bar{x})(y_i-\bar{y})}{\sqrt{\sum_{i=1}^{n}(x_i-\bar{x})^2}\sqrt{\sum_{i=1}^{n}(y_i-\bar{y})^2}},
\label{eq:pearson}
\end{equation}
and Spearman's rank correlation coefficient \citep{spearman1987proof}, defined identically to Eq.~\ref{eq:pearson} but computed on the within-variable ranks of \(x\) and \(y\) rather than on their raw values. Both coefficients range from \(-1\) to \(+1\): the sign indicates the direction of the association (positive: the two quantities tend to increase together; negative: one tends to decrease as the other increases) and the magnitude indicates its strength, with 0 indicating no detectable linear (Pearson) or monotonic (Spearman) relationship. Pearson's \(r\) is sensitive specifically to linear covariation and can be strongly distorted by outliers or by curvature in the relationship; Spearman's \(\rho\) instead only requires that one variable tends to increase (or decrease) monotonically as the other does, without assuming any particular functional shape, making it more robust to the non-Gaussian, often heavy-tailed distributions typical of LEC terms and of the composite-derived descriptors. For this reason, \(\rho\) was adopted as the primary interpretive metric, with \(r\) retained as a linear-association reference.

In practical terms: if larger values of west-sector cold-air advection are systematically associated with larger values of eddy available potential energy (\(A_E\)), the correlation between the two will be positive; if instead stronger advection on that sector tends to coincide with weaker \(A_E\), the correlation will be negative. A pair with \(|\rho|\) close to zero indicates that knowing one quantity provides little information about the other across the cyclone population, whereas \(|\rho|\) approaching 1 indicates a tight, near-deterministic monotonic relationship. Because Pearson's \(r\) and Spearman's \(\rho\) are sensitive to different aspects of the relationship, we treat a substantial divergence between them (\(|\,|r|-|\rho|\,| > 0.10\)) as a diagnostic flag for a non-linear-but-monotonic relationship — for example, an association that strengthens sharply only past some threshold in the predictor, a pattern physically plausible for processes such as frontogenesis, which can intensify non-linearly once a threshold in the thermal gradient is reached.

With \(n=2733\) cyclones entering this analysis, essentially all displayed correlations (\(|r|\ge0.20\)) are statistically significant at conventional thresholds, so — consistent with the interpretive protocol of Sect.~\ref{sec:stats-pvalue} — the correlation results in this study are interpreted primarily on the basis of the magnitude of \(r\) and \(\rho\), the agreement between the two, and physical consistency with the composite fields of Sect.~\ref{sec:composites}, rather than on statistical significance alone. Two further caveats apply to this analysis. First, correlations are computed from contemporaneous, storm-centered, spatially aggregated fields, and therefore diagnose statistical covariability across the cyclone population rather than a mechanistic or causal relationship for any individual storm. Second, the scalar descriptors derived from a given field are not mutually independent (e.g., a border mean and the sector mean on the same side of the cyclone share information), which reduces the number of effectively independent tests being performed; nominal significance levels should therefore be interpreted conservatively, in a field-significance sense \citep{livezey1983statistical,bretherton1999effective,wilks2016stippling}.

\subsubsection{Assessing long-term trends in Energy-Pattern occurrence}
\label{sec:stats-trend}

To determine whether the interannual variability in EP occurrence (Fig.~\ref{fig:6_ep_intensity_seasonality_trends}c) reflects a genuine monotonic trend rather than sampling variability over the 42-year record (1979--2020), we tested the annual cyclone counts of each EP (one value per year, \(n=42\)) using the Mann--Kendall (MK) test \citep{mann1945nonparametric,kendall1948rank}. The MK statistic is
\begin{equation}
S = \sum_{i=1}^{n-1}\sum_{j=i+1}^{n} \mathrm{sign}(x_j - x_i),
\label{eq:mk}
\end{equation}
which simply counts, over every pair of years, whether the later year's count exceeds the earlier one (\(+1\)), is lower (\(-1\)), or is tied (0), and sums the result. Under the null hypothesis of no trend, \(S\) is expected to be near zero and approximately normally distributed for moderate-to-large \(n\); the standardized statistic \(Z = (S-1)/\sqrt{\mathrm{Var}(S)}\) (with the appropriate sign and tie corrections) is compared against the standard normal distribution to obtain a \emph{p}-value. Because the test only asks whether increases outnumber decreases (or vice versa) across all pairs, it does not assume a linear trend, a particular residual distribution, or normality — properties well suited to a 42-point count series that need not follow a simple parametric form.

Trend magnitude was estimated using the Theil--Sen slope \citep{sen1968estimates}, computed as the median of the pairwise slopes \((y_j-y_i)/(x_j-x_i)\) over all \(i<j\):
\begin{equation}
\hat{\beta} = \mathrm{median}\left\{\frac{y_j - y_i}{x_j - x_i} : i < j\right\}.
\label{eq:theilsen}
\end{equation}
Because it is based on the median rather than the mean of pairwise slopes, \(\hat{\beta}\) is markedly less sensitive than an ordinary-least-squares slope to a single anomalous year (e.g., an unusually active or quiet season), and its 95\% confidence interval was estimated from the distribution of pairwise slopes. Trend magnitude and its confidence interval are reported in cyclones~yr\(^{-1}\).

The MK test assumes independent observations; if the annual counts are serially correlated (for example, due to multi-year modes of climate variability), the effective sample size is smaller than \(n=42\) and the nominal MK \emph{p}-value would be anticonservative (too easily significant). We therefore tested for residual autocorrelation prior to reporting MK significance, using the Ljung--Box portmanteau statistic \citep{ljung1978measure},
\begin{equation}
Q(h) = n(n+2)\sum_{k=1}^{h} \frac{\hat{\rho}_k^{\,2}}{n-k},
\label{eq:ljungbox}
\end{equation}
which aggregates the squared autocorrelation \(\hat{\rho}_k\) at lags 1 through \(h\) into a single statistic, tested against a \(\chi^2\) distribution; a significant \(Q(h)\) indicates that the residual series retains structure inconsistent with white noise. This was applied to the residuals obtained after removing the Theil--Sen trend from each EP's annual-count series, using \(h=\min(10, n-1)=10\) lags. None of the three EPs showed significant residual autocorrelation at \(\alpha=0.05\) (EP1: \(p=0.16\); EP2: \(p=0.13\); EP3: \(p=0.05\), marginally above the threshold), so we report trend significance from the original (unmodified) MK test for all EPs, rather than applying a variance correction that is not required by the data and would otherwise reduce statistical power unnecessarily. As a robustness check, we additionally computed and archived MK variants explicitly designed to account for serial correlation — the Hamed--Rao variance correction, the Yue--Wang effective-sample-size correction, pre-whitening, and trend-free pre-whitening \citep{KulkarniVonStorch1995,hamed1998modified,YueEtAl2002_TFPW,YueWang2004} — all of which yielded the same qualitative conclusions regarding which EPs show a statistically significant trend (not shown). In Fig.~\ref{fig:6_ep_intensity_seasonality_trends}c, solid lines denote EPs with a significant monotonic trend (\(p<0.05\)) and dashed lines denote non-significant trends, with the sign of the Theil--Sen slope indicating the trend direction in either case.

Taken together, the sequential decision procedure, effect sizes, multiple-comparison corrections, correlation diagnostics, and trend tests described above form a single, consistent statistical framework applied throughout the Results: statistical significance is used to flag candidate differences, associations, or trends, while effect size, agreement between complementary metrics (e.g., Pearson vs.\ Spearman, or the different MK variants), and consistency with the independently derived composite dynamical fields (Sect.~\ref{sec:composites}) determine which of these candidates are interpreted as physically robust features of the energy patterns.


\section{Results and Discussion}\label{sec:results}

\subsection{General energetic behavior of Southwestern Atlantic cyclones}

Figure \ref{fig:3_phase_density_2x2} contains the conversion and imports LPS diagrams for all cyclones considered in the current study, averaged for each cyclone lifecyle phase. On then, the bins indicate mean $G_E$ values and the contours, cyclone density. The cyclone density always increase towards the graph origin and decreases outwards, with the regions of higher density changing for each develoment phase. This indicates that the cyclones behavior represent a continuum, a feature ubiquoustly noted by all previous studies attempting an objective cyclone classification \citet[e.g.][]{hart2003cyclone,yanase2014parameter,graf2017objective,dekker2018characteristics,corner2025classification}.

For the conversion LPS, there is a tendency for systems to concentrate in the upper-right quadrant across all phases (Fig.~\ref{fig:3_phase_density_2x2}a,c,e,g). This indicates a shared importance of baroclinic and barotropic processes for eddy development, as previously noted by \citet{de2025lorenz}. This concentration strengthens from the incipient to the intensification phase and then gradually weakens. The $G_E$ term follows a similar evolution: although values are averaged within each $3 \times 3$ bin, an increase in the relative contribution of $G_E$ from incipient to intensification can be identified, followed by a decrease afterward. Also, $G_E$ increase approximately with increasing $C_A$. A smaller fraction of cases falls in regimes in which barotropic processes alone govern cyclone development (lower-left quadrants) or in which neither baroclinic nor barotropic conversions favor eddy growth (lower-right quadrants). For cyclones in the lower-left quadrant, $G_E$ contributions are generally weaker. Cyclones in the lower-right quadrant (which are rare, accounting for roughly 5\%) are either driven by $G_E$ or by energy imports.

The ``import'' LPS exhibits greater dispersion than the ``conversion'' LPS. Although cyclone behavior spans a continuum, some phase-dependent preferences emerge. In the incipient phase (Fig.~\ref{fig:3_phase_density_2x2}b), there is a concentration in the right half-plane, indicating a preference for imports of $A_E$ during this phase, whereas $K_E$ may be either exported or imported. During intensification (Fig.~\ref{fig:3_phase_density_2x2}d), the relative importance of the left half-plane increases, particularly the upper-left quadrant, indicating a higher tendency to export $A_E$ while importing $K_E$. During maturity (Fig.~\ref{fig:3_phase_density_2x2}f), there is a stronger concentration in the lower-left quadrant, indicating exports of both forms of eddy energy. Lastly, during decay (Fig.~\ref{fig:3_phase_density_2x2}h), cyclones concentrate in either the lower-left or upper-right quadrants, suggesting that, in some cases, energy export contributes to cyclone dissipation, whereas in others, energy import may help maintain the system despite frictional processes \citep{de2025lorenz} and/or enhanced conversion back to the mean flow.


\begin{figure}[ht!]
    \centering
    \includegraphics[width=1\linewidth]{3_phase_density_2x2.png}
    \caption{Phase-dependent density of cyclone energetics in the Lorenz Phase Space (LPS). Conversion LPS panels show $C_K$ (x-axis) versus $C_A$ (y-axis) for the incipient (a), intensification (c), mature (e), and decay (g) phases. Import LPS panels show $BA_E$ (x-axis) versus $BK_E$ (y-axis) for the corresponding phases (b,d,f,h). Colored bins indicate the mean $G_E$ within each bin, and black contours depict the relative density of samples. Percentages in the corners denote the fraction of samples in each quadrant. The colored map was obtained by discretizing the 2‑D phase space into a 20×20 regular grid and computing the arithmetic mean of $G_E$ in each cell; cells with fewer than 3 observations were masked to reduce statistical noise. Density contours are derived from the per‑cell counts smoothed with a Gaussian filter.}
    \label{fig:3_phase_density_2x2}
\end{figure}

\subsection{Energy patterns and key energy pathways for cyclone development}
\label{sec:energy_patterns}

The clustering analysis yielded three centroids defining the energy patterns (EPs) of cyclones with genesis in the southwestern Atlantic. Pairwise effect sizes based on the rank-biserial correlation coefficient indicate a clear energetic separation among the three EPs, with non-negligible differences for most LEC terms in all pairwise contrasts (Figure~\ref{fig:S2_pairwise_effectsize_lec_terms}). Notably, most EP1--EP2 comparisons already show meaningful effect sizes, although the strongest separations are generally associated with contrasts involving EP3. Non-negligible differences are also evident for several terms not directly included in the clustering procedure, supporting the robustness of the segmentation and indicating that the EPs capture broader and physically consistent distinctions in cyclone energetics.

Figure~\ref{fig:4_lps_combined}a shows the conversion LPS for these EPs, indicating that, for all EPs, the energy cycle is largely modulated by barotropic and baroclinic conversions that feed the eddy. For EP3, these contributions are more modest, including the contribution from $G_E$. For EP2, systems begin the incipient phase with moderate conversions, which increase nearly proportionally during intensification, return approximately to their initial values, and then decrease toward near-neutral conditions. EP1 also begins with modest conversions, but exhibits large increases during intensification and further increases during the mature phase, particularly in barotropic conversion. EP1 also shows relatively large conversion values during decay. $G_E$ is highest during intensification for both EP1 and EP2, but $K_E$ is larger for EP2 during intensification, suggesting that EP1 systems tend to intensify more slowly.

The $C_A$ and $C_K$ terms are now assessed only during the intensification phase, when these terms reach their largest magnitudes. In terms of $C_A$ magnitude, EP1 and EP2 are statistically similar (Fig.~\ref{fig:S2_pairwise_effectsize_lec_terms}), with values peaking near the surface, remaining relatively large from low to mid-levels, and increasing again between 600 and 300\,hPa (Fig.~\ref{fig:S3_vertical_levels}a). In contrast, EP1 exhibits significantly larger $C_K$ (Figure~\ref{fig:4_lps_combined}a); therefore, we further investigate the characteristics of this term for EP1 systems. For EP1, $C_K$ reaches its largest magnitudes between 400 and 250\,hPa (Fig.~\ref{fig:5_ck_subterms_vertical_profiles}a). This vertical structure is broadly mirrored by EP2 cyclones, whereas EP3 cyclones show consistently low $C_K$ values throughout the atmosphere (Fig.~\ref{fig:S3_vertical_levels}b).

As shown in Fig.~\ref{fig:5_ck_subterms_vertical_profiles}b, the $C_K$ subterm with the largest median magnitude is term E, which is related to the vertical flux of eddy meridional momentum coupled with the vertical shear of the mean meridional wind. This is followed by term B, associated with the meridional flux of eddy kinetic energy and the meridional gradient of the mean meridional wind, and term A, which represents the primary barotropic instability mechanism. Terms D and C exhibit low median magnitudes and are mostly positive, indicating that they contribute to energy transfer from the eddy to the mean flow. Term D is associated with the vertical flux of eddy zonal momentum coupled with the vertical shear of the mean zonal wind, whereas term C represents eddy zonal kinetic energy modulated by the mean meridional flow. Although term E has a larger median magnitude than term B, the latter exhibits greater variability and is therefore more intense in some systems. Among EP1 cyclones, term E is dominant in 43\% of the cases, followed by term B in 38\% and term A in 19\%.

\begin{figure}[ht!]
    \centering
    \includegraphics[width=1\linewidth]{5_ck_subterms_vertical_profiles.png}
    \caption{$C_K$ subterm characteristics for EP1 cyclones during the intensification phase.  (a) Vertical distribution of $C_K$ as a function of pressure level. Boxplots show the distribution across cyclones, with the red line indicating the median and whiskers representing the interquartile range.(b) Domain-integrated $C_K$ subterms (A–E) across cyclones. The dashed vertical line marks zero, separating energy transfer from the mean flow to the eddy (negative values) from transfer from the eddy to the mean flow (positive values).}
    \label{fig:5_ck_subterms_vertical_profiles}
\end{figure}


Figure~\ref{fig:4_lps_combined}b shows the import LPS. EP3 initially exhibits modest imports of $A_E$ that decrease over time, while $K_E$ remains close to neutral; by decay, both terms are nearly neutral. EP2 shows a similar behavior, but with modest imports of $K_E$; both imports decrease over the life cycle, becoming nearly neutral by decay. EP1, in contrast, starts with small imports of $A_E$ and near-neutral $K_E$; while $K_E$ imports increase slightly during intensification, the system transitions to exporting $A_E$. During the mature and decay phases, EP1 exports large amounts of both forms of eddy energy. This suggests that the large conversion values attained during these phases counteract these exports (and dissipation; \citealp{de2025lorenz}), keeping the system active until it eventually dissipates.

\begin{figure}[ht!]
    \centering
    \includegraphics[width=1\linewidth]{4_lps_combined.png}
    \caption{Centroids of the three energy patterns (EP1--EP3) in the Lorenz Phase Space (LPS), shown for phase-mean energetics across the cyclone life cycle. (a) Conversion LPS: $C_K$ (x-axis) versus $C_A$ (y-axis). (b) Import LPS: $BA_E$ (x-axis) versus $BK_E$ (y-axis). Markers represent life-cycle phase means (incipient, intensification, mature, and decay) connected in chronological order; marker size scales with $K_E$ and marker color indicates $G_E$. Line thickness denotes the EP group. Letters A and Z denote the beginning and end of each trajectory, respectively.}
    \label{fig:4_lps_combined}
\end{figure}

Figure~\ref{fig:6_ep_intensity_seasonality_trends} summarizes the overall characteristics of cyclones within each EP. Overall, EP3 includes the weakest systems in the dataset (Fig.~\ref{fig:6_ep_intensity_seasonality_trends}a), consistent with its comparatively less active energy cycle, and is also the most frequent group (Fig.~\ref{fig:6_ep_intensity_seasonality_trends}c). These systems exhibit pronounced seasonality, with higher genesis frequency during DJF and SON (Fig.~\ref{fig:6_ep_intensity_seasonality_trends}b). In contrast, EP1 comprises the strongest systems, is also the least frequent, and shows a clear peak during JJA. EP2 displays intermediate intensity and a weaker seasonal modulation, with relatively higher frequencies during JJA and SON.

\begin{figure}[ht!]
    \centering
    \includegraphics[width=1\linewidth]{6_ep_intensity_seasonality_trends.png}
    \caption{Overall characteristics of the energy-pattern groups (EP1--EP3). (a) Distribution of cyclone maximum intensity, quantified as the peak value of $-\zeta_{850}$ (1 $\times$ 10$^{-5}$\,s$^{-1}$); horizontal bars and labels indicate the maximum, mean, median and mean for each EP. (b) Seasonal distribution of cyclone genesis frequency (DJF, MAM, JJA, SON) for each EP. (c) Interannual variability of annual cyclone counts (1979--2020) for each EP; solid (dashed) lines denote significant (non-significant) monotonic trends at $p<0.05$.}
    \label{fig:6_ep_intensity_seasonality_trends}
\end{figure}

The relative contribution of each genesis region for each EP is shown in Fig.~\ref{fig:7_ep_genesis_density_kde}. Figure~\ref{fig:7_ep_genesis_density_kde}a shows the genesis density for all systems, indicating the highest frequency in the ARG region, followed by LA-PLATA and then SE-BR, as previously reported \citep{hoskins2005new,gramcianinov2019properties,crespo2021potential,couto2024new}. Figures~\ref{fig:7_ep_genesis_density_kde}b, c, and d present the difference in genesis density when comparing the density of each EP individually with the density of all cyclones. For this, both densities were normalized using min-max normalization. The normalized genesis density for each EP was then subtracted from the normalized density of all systems (i.e., relative anomaly = norm\_EP - norm\_all). Relative to the all-cyclone distribution, EP1 shows a higher relative frequency of genesis in SE-BR and in an intermediate region between ARG and LA-PLATA (Fig.~\ref{fig:7_ep_genesis_density_kde}b). EP2 exhibits a higher relative occurrence of genesis in LA-PLATA (Fig.~\ref{fig:7_ep_genesis_density_kde}c), whereas EP3 shows a smaller relative frequency across all regions, with an enhanced relative frequency south of ARG (Fig.~\ref{fig:7_ep_genesis_density_kde}d).

\begin{figure}[ht!]
    \centering
    \includegraphics[width=1\linewidth]{7_ep_genesis_density_kde.png}
    \caption{Kernel density estimates of cyclone genesis density (cyclones per $10^6$ km$^2$ per year) for all cyclones (a) and relative frequency bias for each energy-pattern group: (b) EP1, (c) EP2, and (d) EP3. Shading and contours on (a) represent genesis density (cyclones per $10^6$ km$^2$ per year). In panels (b--d), the genesis density for each EP was first normalized using min-max normalization and then subtracted from the normalized density for all systems, resulting in the relative anomaly of genesis density for each EP.}
    \label{fig:7_ep_genesis_density_kde}
\end{figure}

Overall, these results indicate that EP3 comprises weaker, day-to-day systems, with genesis predominantly in ARG. These cyclones have a less active energy cycle and occur preferentially during the austral warm season. Although they exhibit substantial interannual variability, no statistically significant trend is detected in their occurrence. In contrast, EP2 represents relatively intense systems, with genesis primarily in LA-PLATA and ARG, weak seasonality, and energetics characterized by modest baroclinic conversions and diabatic contributions, together with especially large imports of $A_E$ and $K_E$ during early development. EP2 is the only EP exhibiting a statistically significant trend, with an overall increase.

Finally, EP1, while being the least frequent and most active during the cold season, shows the strongest conversions throughout the life cycle, particularly barotropic conversion, with relevant contributions from diabatic processes. However, these eddy-energy sources are partly offset by energy export, which becomes stronger during the mature and decay phases. EP1 does not show a marked interannual variability or a statistically significant trend.


\subsection{Cyclone structure}\label{sec:composites}

Figure \ref{fig:8_dynamical_composites_epall_relative} presents the dynamical composites for EP1, EP2, and EP3. The first row (\ref{fig:8_dynamical_composites_epall_relative}a--c) shows upper-level PV ($PV'_{200}$), lower- to mid-tropospheric baroclinicity (represented by $EGR_{850-500}'$), and sea-level pressure. For PV and EGR, anomalies are shown with respect to the total cyclone composite (EPALL), whereas for SLP the total composite is shown. For all EPs, the SLP field exhibits a clear cyclonic signature, marked by closed isobars. For both EP2 and EP3 there is a ridge downstream of the surface cyclone. However, for EP1 the ridge seems to be upstream of the surface cyclone, a feature which will be relevant for future analysis. 

The upper-level PV anomaly patterns differ substantially among the EPs, indicating distinct large-scale dynamical configurations. For EP1, the $PV'_{200}$ field exhibits an elongated structure that resembles a PV-streamer-like signature, which is consistent with an equatorward intrusion of high stratospheric PV \citep[e.g.,][]{hoskins1985use,hu2004variations,wernli2007identification}. Because this is an EPALL-relative anomaly field, this pattern should be interpreted as streamer-like rather than as a definitive PV streamer. In this case, there is no equally compact or collocated cyclonic signature in the lower-tropospheric PV anomaly ($PV'_{850}$; \ref{fig:8_dynamical_composites_epall_relative}d), suggesting weaker vertical coupling between the upper- and lower-level anomalies. 

In contrast, EP2 displays a more compact $PV'_{200}$ structure, with a cutoff-like or strongly wrapped trough--ridge dipole, which appears more clearly coupled to the lower-level PV anomaly ($PV'_{850}$; Fig. \ref{fig:8_dynamical_composites_epall_relative}e). This configuration is consistent with a more favorable vertical coupling between upper- and lower-tropospheric PV anomalies, which may support more efficient cyclone intensification than in the more elongated EP1 configuration \citep{davis1991potential,huo1999interaction,schemm2019efficiency}. EP3 presents a upper and low-level $PV'$ signal inverted than EP2, indicating a weaker  dynamical coupling and thus, weaker development. 


\begin{figure}[ht!]
    \centering
    \includegraphics[width=1\linewidth]{8_dynamical_composites_epall_relative.png}
    \caption{Storm-centered composite anomalies relative to the full cyclone population (EPALL) at the central
    intensification timestep. Each column represents one energy pattern: EP1 ($n = 332$), EP2 ($n = 775$), and EP3 ($n = 1{,}623$). Anomalies are computed as the difference between the energy-pattern composite and the EPALL composite ($n = 3{,}820$ cyclones). All fields are displayed on a storm-centered $30^\circ \times 30^\circ$ grid; the black star marks
    the cyclone center and the dashed rectangle delimits the $15^\circ \times 15^\circ$ domain used for LEC computations. Top row (a--c): Shading shows the anomaly of potential vorticity at 200~hPa ($\mathrm{PV}_{200}'$, contours show the Eady growth rate anomaly ($\mathrm{EGR'_{850-500}}$): positive contours in solid red at $+0.03$, $+0.06$, $+0.09$, and $+0.12$~day$^{-1}$; negative contours in dashed blue at $-0.12$, $-0.09$, $-0.06$, and $-0.03$~day$^{-1}$. Arrows show anomaly wind vectors at 250~hPa ($\mathbf{V'}_{250}$). Thin semi-transparent black contours show total sea-level pressure (interval $2$~hPa). Middle row (d--f): Shading shows the anomaly of potential vorticity at 850~hPa ($\mathrm{PV'}_{850}$. Contours show the anomaly of 850~hPa temperature advection ($\mathrm{AdvT'}_{850}$): warm advection in solid red and cold advection in dashed blue, with an interval of  $0.05$~K~h$^{-1}$ and a range of $\pm\,0.20$~K~h$^{-1}$. Arrows show anomaly wind vectors at 850~hPa ($\mathbf{V'}_{850}$). Bottom row (g--i): Shading shows the anomaly of ageostrophic geopotential flux convergence at 250~hPa ($\mathrm{AFC'}_{250})$. Contours show the anomaly of kinetic energy advection at 250~hPa ($\mathrm{AdvKE'}_{250}$): positive in solid red at $+0.005$ and $+0.010$~m$^2$~s$^{-3}$; negative in dashed blue at
    $-0.010$ and $-0.005$~m$^2$~s$^{-3}$. Arrows show anomaly wind vectors at 250~hPa ($\mathbf{V'}_{250}$). Hatched regions indicate where the RK criterion undergoes a meridional sign reversal in the \emph{total} composite field (not the EPALL-relative anomaly).} \label{fig:8_dynamical_composites_epall_relative}
\end{figure}

The $EGR_{850-500}'$ fields are also markedly distinct among the EPs. For EP1, enhanced baroclinicity is found around the cyclone center and extends downstream of it. In contrast, for EP2, the baroclinic zone is elongated in a NW--SE direction, similarly to the configuration described by \citet{crespo2021potential} for cyclones with genesis in the Uruguay (LA-PLATA) region. This suggests that these systems develop within a strongly baroclinic environment embedded in the South Atlantic storm-track region \citep{hoskins1990existence,chang2002storm,hoskins2005new}. The observed pattern further suggests differences in the relative positioning between the cyclone core and the surrounding low-level baroclinic reservoir in EP1 and EP2. In both cases, however, the upper-level PV anomalies appear to interact with a favorable baroclinic environment, supporting ascent and cyclonic intensification \citep{davis1991potential,huo1999interaction,besson2021global,crespo2021potential}. Lastly, EP3 cyclones exhibit a weaker relative baroclinic signature in the same region where EP2 shows its maximum $EGR_{850-500}$ anomaly.
 

Figure \ref{fig:8_dynamical_composites_epall_relative}d-f presents low-level PV ($PV'_{850}$), temperature advection ($AdvT'_{850}$) and wind vectors at 850 hPa - all anomalies in respect to EPALL. In respect to $PV'_{850}$, as previously discussed, EP1 presents a less spatially coherent field, while EP2 presents a cyclonic PV anomaly coupled to the upper-level disturbance. EP3 presents a weak anticyclonic anomaly near the cyclonic center, indicative of a weaker cyclonic development.

EP2 presents a clear cold--warm dipole embedded in a cyclonic circulation, consistent with the canonical frontal structure of extratropical cyclones \citep[e.g.][]{dacre2025kinematic}. In this sense, the pattern agrees with South Atlantic cyclone composites associated with organized frontal development \citep{gramcianinov2019properties}. Warm advection in the cyclone warm sector is, according to quasigeostrophic theory, associated with ascent along and ahead of the warm front \citep{davies2015quasigeostrophic}. However, the resulting positive low-level PV anomalies are more directly linked to the diabatic processes accompanying this ascent, particularly condensation and latent heating, rather than to thermal advection alone \citep{hoskins1985use,huo1999interaction,attinger2021systematic}.

EP1, in turn, is dominated by relative cold-advection anomalies, together with an anticyclonic circulation southwest of the surface cyclone, suggestive of a post-frontal environment. In combination with the streamer-like upper-level PV structure, this configuration is consistent with a post-frontal or secondary-development setting, although it is not by itself sufficient to diagnose secondary cyclogenesis. A similar South Atlantic configuration was described by \citet{iwabe2009event}, while \citet{gramcianinov2019properties} also suggested that some cyclones forming near southeastern Brazil and the southeastern South Atlantic are associated with secondary development. Lastly, EP3 exhibits an anomaly pattern opposite to EP2, indicating a weaker or displaced canonical frontal thermal-advection dipole relative to the EPALL composite. In this sense, EP3 does not appear to display the same degree of organized frontal thermal contrast found in EP2.

Figure \ref{fig:8_dynamical_composites_epall_relative}g--i presents anomalies of upper-level ageostrophic geopotential flux convergence (AFC$_{250}'$) and kinetic energy advection ($\mathrm{AdvKE'}_{250}$), together with hatching indicating regions where the meridional gradient of absolute vorticity changes sign, that is, where the Rayleigh--Kuo (RK) necessary condition for barotropic instability is satisfied. The RK hatching is shown for the total EP composite rather than for the anomalous fields. Since the Rayleigh--Kuo criterion provides only a necessary, not sufficient, condition for barotropic instability, the hatching should be interpreted as background dynamical context rather than as direct evidence of barotropic growth.

The AFC$_{250}'$ field highlights important distinctions between EP1 and EP2. EP2 presents a dipolar pattern, with relatively intense positive anomalies on the western side of the cyclone and relatively intense negative anomalies on its eastern side. This indicates enhanced local upper-level eddy kinetic energy accumulation to the west of the cyclone core and depletion to the east, relative to EPALL, a configuration consistent with an organized downstream-dispersive upper-level energetics pattern during cyclone intensification \citep{orlanski1993ageostrophic,orlanski1995stages,piva2010energetics,pinheiro2022contributions}. The upper-level wind anomalies are consistent with this interpretation, as EP2 exhibits a more coherent cyclonic-anomalous circulation around the cyclone core, with stronger meridional flow anomalies on the flanks of the system, suggesting a more sharply defined shear and deformation environment. The $\mathrm{AdvKE'}_{250}$ anomalies are broadly consistent with this interpretation, since EP2 also exhibits a dipolar structure, suggesting a more organized redistribution of kinetic energy along the upper-level flow, with positive anomalies downstream of the surface cyclone.

In contrast, EP1 exhibits a broader positive AFC$_{250}'$ anomaly over and downstream of the cyclone center, while negative anomalies are found mainly upstream and to the southwest. Relative to EPALL, this suggests a less canonical downstream-dispersive AFC pattern than in EP2, with comparatively greater local upper-level energy accumulation in the vicinity of and downstream of the cyclone core. In this sense, the downstream development signature appears less clearly organized in EP1, suggesting that other energetic pathways may play a relatively larger role in sustaining cyclone growth. For EP1, the surface cyclone appears anomalously closer to the upper-level trough axis, that is, downstream of an upstream ridge and less horizontally displaced from the upper-level PV anomaly than in EP2. This suggests a more trough-controlled upper-level configuration, consistent with upper-level-induced cyclogenesis or with systems developing beneath a pre-existing upper-level trough, rather than with the more canonical downstream-dispersive pattern inferred for EP2 \citep{schultz2019extratropical,sinclair2000classification,graf2017objective}.
 The negative $\mathrm{AdvKE'}_{250}$ anomalies reinforce the interpretation that kinetic energy redistribution in EP1 differs from that of EP2, although this term should be viewed primarily as a redistribution term rather than as a direct source or sink of cyclone intensification.

EP3, in turn, is characterized by predominantly negative AFC$_{250}'$ anomalies around the cyclone core, indicating weaker upper-level energetic support relative to EPALL. Its upper-level wind anomalies are weaker and more diffuse than in EP2, which is consistent with the less organized AFC pattern. Finally, all EPs exhibit RK sign reversals close to the cyclone center, suggesting that the necessary condition for barotropic instability is locally satisfied in all of them. This supports the interpretation that barotropic processes may be present in the three EPs, as a more localized phenomenon near the upper-level shear zone \citep[e.g.][]{mak1989local,wernli1998effect,cai2004local,riviere2006role}.

\subsection{Relationship between energy terms and dynamical fields}\label{sec:correlation}

To assess how LEC terms covary with the dynamical fields shown in Sect.~\ref{sec:composites}, we computed Pearson linear correlations ($r$) between each canonical LEC term and each scalar feature derived from each field, and used Spearman rank correlation ($\rho$) as a monotonic-relationship check. Figure~\ref{fig:9_pearson_epall_by_field_type} displays only pairs with $|r| \geq 0.20$.  Although this threshold is usually considered to indicate a weak relationship, the large variability in the cyclone dataset and the complexity of LEC terms (which often involve multiple contributions, including departures from the zonal mean, deviations from the domain mean, and different forms of spatial averaging) make these
associations noteworthy.  The $\mathrm{PV}_{850}$ field yielded only 7 of 91 possible pairs (7.7\%) above the display threshold, and therefore is not accessed. Similarly, the scalar fields derived from RK do not present any meaningful physical meaning and are also not accessed.

In this subset ($126$ pairs), Pearson and Spearman signs agree in all cases (100\%), and only 5 pairs show $|\,|r|-|\rho|\,|>0.10$. Given $n=2733$, all displayed correlations are highly significant statistically, but interpretation below is based primarily on effect size and physical consistency rather than on $p$-values alone. A caveat is that these are contemporaneous, storm-centered, spatially aggregated features; therefore they diagnose covariability, not direction of causality. In addition, dependence among features (e.g., shared information between border, sector, and contrast metrics) reduces the number of effectively independent tests, so nominal significance levels should be interpreted conservatively in a field-significance sense \citep{livezey1983statistical,bretherton1999effective,wilks2016stippling}.

\begin{figure}[ht!]
    \centering
    \includegraphics[width=1\linewidth]{9_pearson_epall_by_field_type.png}
    \caption{Pearson correlation matrix between scalar variables derived from the dynamical fields used in the composite analysis and the LEC terms. Only correlations with $|r| > 0.2$ are shown. Hatches indicate combinations where $|\,|r| - |\rho|\,| > 0.10$}
    \label{fig:9_pearson_epall_by_field_type}
\end{figure}

For $\mathrm{AdvT}_{850}$ (Fig.~\ref{fig:9_pearson_epall_by_field_type}a), the strongest links involve APE-related terms. For $C_A$, the largest magnitudes are in sector W $(r=-0.46,\,\rho=-0.55)$, sector N $(r=-0.45,\,\rho=-0.56)$, and domain mean $(r=-0.44,\,\rho=-0.50)$. For $A_E$, the largest values are in sector N $(r=-0.56,\,\rho=-0.64)$ and border N $(r=-0.54,\,\rho=-0.56)$. These signs are consistent with a stronger cold-advection signal on the cold side of the cyclone, inducing a steeper meridional temperature gradient, being associated with enhanced baroclinic conversion and larger eddy APE, in line with quasigeostrophic/baroclinic-growth theories \citep{lorenz1955available,davies2015quasigeostrophic}. A notable divergence between Pearson and Spearman for $C_A \times \mathrm{sector\_N}$ ($|r| = 0.44$ vs $|\rho| = 0.55$) suggests that the $C_A$–temperature advection relationship has a moderately non-linear component, physically consistent with the intensification of frontal zones on the equatorward side of extratropical cyclones, where frontogenesis, thermal advection, and warm--moist pre-frontal transport may accompany or favor rapid cyclone deepening \citep{dacre2025kinematic}.

For $BA_e$, $\mathrm{AdvT}_{850}$ correlations show a coherent east/west dipole: border E $(r=+0.28,\,\rho=+0.29)$, sector E $(r=+0.28,\,\rho=+0.30)$, and E--W contrast $(r=+0.33,\,\rho=+0.31)$ are positive, whereas border W $(r=-0.24,\,\rho=-0.20)$ and sector W $(r=-0.37,\,\rho=-0.28)$ are negative. The physical interpretation is straightforward: warm (cold) influx on the eastern side increases (decreases) $BA_E$, whereas warm (cold) advection on the western side decreases (increases) $BA_E$. Weak but significant correlations involving $K_E$ ($|r| = 0.22$–$0.35$) are likely indirect: $C_A$ and $BA_e$ increase the $A_E$ reservoir, which is subsequently converted into $K_E$ through $C_E$.


For $\mathrm{AFC}_{250}$ (Fig.~\ref{fig:9_pearson_epall_by_field_type}b), the dominant signal is the amplitude metric $|\mathrm{AFC}|$ (domain absolute mean), most strongly correlated with $K_E$ $(r=+0.53,\,\rho=+0.57)$ and $A_E$ $(r=+0.51,\,\rho=+0.56)$. The signed domain-mean AFC is positively associated with $C_A$ $(r=+0.43,\,\rho=+0.38)$ and negatively with $C_K$ $(r=-0.22,\,\rho=-0.23)$. This indicates that stronger upper-level ageostrophic geopotential flux convergence is associated with larger eddy energy reservoirs, broadly consistent with downstream-development energetics \citep{orlanski1991life,orlanski1993ageostrophic}, but should be interpreted as population-level co-variability rather than event-scale mechanism proof. The large absolute-value correlation, rather than the signed mean, implies that the key signal is the intensity of the upper-level energy redistribution, regardless of its direction, rather than a systematic bias toward convergence or divergence. This is in accordance with the AFC composites (Figure \ref{fig:8_dynamical_composites_epall_relative}g--i), as EP1 and EP2, which represent the most intense systems, present higher energy divergence and convergence, respectively.

Regarding the signed domain-mean AFC, $C_A$ shows $r = +0.43$ and $C_k$ shows $r = -0.22$. The positive $C_A$–AFC relation suggests that, on average, cyclones with more organized upper-level downstream energy redistribution tend to have stronger baroclinic and barotropic conversions. For $C_A$, positive sigifficant correlations are found also for the magnitude of domain mean, central value and all sectors. Meanwhile, for $C_K$ it is found only for the eastern flank. This suggests that barotropically dominant cyclones (EP1) are associated with a broader, more symmetric AFC pattern, rather than the predominantly eastward-dispersing AFC signature typical of classically
baroclinic downstream development \citep{orlanski1991life}, as noted in the composite analysis for EP2.

The spatial pattern of AFC correlations with $BA_E$ and $BK_E$ are also noteworthy, with the correlation sign reversed. For $BA_E$, border/sector W are positive [border W: $(r=+0.38,\,\rho=+0.37)$; sector W: $(r=+0.20,\,\rho=+0.19)$], while border E and E--W contrast are negative [border E: $(r=-0.29,\,\rho=-0.27)$; contrast E--W: $(r=-0.41,\,\rho=-0.40)$]. This supports an interpretation in which west-enhanced AFC configurations covary with stronger net $A_e$ import, while the reverse is apparent for $BK_E$.

For $\mathrm{KE\_adv}_{250}$ (Fig.~\ref{fig:9_pearson_epall_by_field_type}c), the strongest individual pair is $A_E\times$sector W $(r=-0.44,\,\rho=-0.45)$. The amplitude metric $|\mathrm{KE\_adv}_{250}|$ is positively correlated with $A_E$ $(r=+0.43,\,\rho=+0.48)$, $C_A$ $(r=+0.34,\,\rho=+0.40)$, and $K_E$ $(r=+0.33,\,\rho=+0.40)$. This indicates that the overall upper-level energy redistribution efficiency correlates with the baroclinic energy pathway. Also, the correlation signals for $\mathrm{AdvKE}_{250}$ borders and sectors are largely opposite of those found by AFC. This is in accordance with EP2 (more intense) and EP3 (less intense) composites (Figure \ref{fig:8_dynamical_composites_epall_relative}h,i) displaying dominantly negative and positive $\mathrm{AdvKE'}_{250}$ values on the cyclone western sector, respectively.

For $C_K$, positive associations occur for domain mean and sectors N/S/E [domain mean: $(r=+0.33,\,\rho=+0.37)$; sector N: $(r=+0.26,\,\rho=+0.26)$; sector S: $(r=+0.23,\,\rho=+0.30)$; sector E: $(r=+0.27,\,\rho=+0.29)$]. These positive correlations imply less negative $C_k$ where KE-advection metrics are larger. consistent with cyclones with the strongest barotropic instability (EP1) being characterized by large eddy-energy export or weaker imports (Figure \ref{fig:4_lps_combined}). For $BK_E$, the only correlation above the threshold is with the western border ($(r=+0.37,\,\rho=+0.33)$), suggesting that KE advection from the west directly modulates the lateral $K_E$ import.

For $\mathrm{PV}_{200}$ (Fig.~\ref{fig:9_pearson_epall_by_field_type}d), the largest coefficients in the full figure are tied to E--W contrast: $G_E$ $(r=+0.59,\,\rho=+0.57)$, $A_E$ $(r=+0.57,\,\rho=+0.59)$, and $C_A$ $(r=+0.45,\,\rho=+0.47)$. A positive E--W contrast corresponds to less-negative PV to the east and more-negative PV to the west. This is consistent with a trough-west/ridge-east configuration associated with stronger diabatic generation and baroclinic energetics \citep{hoskins1985use,davis1991potential}.

Consistently, border/sector W PV values are negatively correlated with $C_A$ and $A_E$ [border W: $(r_{C_A}=-0.30,\,\rho=-0.31)$ and $(r_{A_E}=-0.44,\,\rho=-0.47)$; sector W: $(r_{C_A}=-0.24,\,\rho=-0.25)$ and $(r_{A_E}=-0.37,\,\rho=-0.39)$], indicating larger $C_A$ and $A_E$ under more cyclonic (more negative) western PV. By contrast, the isolated relation between border-W PV and $G_E$ is weak $(r=-0.05,\,\rho=-0.01)$, so this quantity should not be interpreted as a robust standalone predictor of $G_E$.

\section{Summary and Conclusions}

- LPS se mostrou uma ferramenta útil para visualização dos principais termos do LEC relacionados ao desenvolvimento dos ciclones, de forma a condensar a informação e apresentar de um modo intuitivo em comparação ao tradicional diagrama do LEC com as quatro caixas e setas de conversão. Não que vá substitui-lo, mas para análise de sistemas transientes é uma boa opção.

-  While conversion LPS shows that the difference among systems is more represented by the instensity of the Ck and Ca conversions, the imports LPS shows distinct pathways.

- EP1 high conversion on the decay phase phase act to maintain the system against energy dissipation \citep{de2025lorenz} and exports.

- EP2: classic baroclinic systems, with strong AFC convergence; EP1: possibly secondary cyclogenesis, with weakened AFC, and strong barotropic conversions. But all systems presented both baroclinic and baroclinic conversions. Ibc was expected, but the dominant role of Ibt is somewhat unexpcted and localized. 

- Importação de AE pode estar relacionado com transporte remoto de calor, enquanto que importação de KE pode ser o downstream development - mais estudos são necessários.

- A analise de correlação indicou que o LEC no framework semi-lagrnagiano é capaz de representar campos dinamicos fisicamente realistas e trazer novos insights sobre o desenvolvimento desses sistemas. 


\backmatter

\section*{Supplementary information}

\renewcommand{\thefigure}{S\arabic{figure}}
\setcounter{figure}{0}

\begin{figure}
    \centering
    \includegraphics[width=1\linewidth]{S1_pca_clustering_validation.png}
    \caption{Principal component analysis (PCA) and selection of the optimal number of $k$-means clusters. (a) PCA scree plot showing the fraction of variance explained by each principal component (bars) and the cumulative explained variance (green line). The horizontal red dashed line marks the 90\% threshold, and the vertical red dashed line indicates the retained dimensionality ($K=15$ PCs), which explains 90.3\% of the cumulative variance. (b) Determination of the optimal number of clusters ($k$) using normalized cluster-validity indices computed in the retained PCA space. Individual curves show the min--max normalized indices (with Davies--Bouldin inverted so that higher values indicate better performance); the black curve is the unweighted mean index. The vertical red dashed line and star highlight the selected solution, $k=3$.}
    \label{fig:S1_pca_clustering_validation}
\end{figure}

\begin{figure}
    \centering
    \includegraphics[width=1\linewidth]{S2_pairwise_effectsize_lec_terms.png}
    \caption{}
    \label{fig:S2_pairwise_effectsize_lec_terms}
\end{figure}

\begin{figure}
    \centering
    \includegraphics[width=1\linewidth]{S3_vertical_levels.png}
    \caption{Vertical distribution of conversion terms for EPs 1-3 ($W m ^{-2}$. (a) Baroclinic conversion ($C_A$) by pressure level; the star marks the level of maximum median $C_A$. (b) Barotropic conversion ($C_K$) by pressure level; the star marks the level of minimum median $C_K$.}
    \label{fig:S3_vertical_levels}
\end{figure}

\begin{figure}
    \centering
    \includegraphics[width=1\linewidth]{S4_pairwise_effectsize_composite_scalars.png}
    \caption{}
    \label{fig:S4_pairwise_effectsize_composite_scalars}
\end{figure}

%TC:ignore
\section*{Declarations}

\subsection*{Funding}

This study was partly financed by the Coordenação de Aperfeiçoamento de Pessoal de Nível Superior – Brasil (CAPES) under Finance Code 001.

\subsection*{Competing interests}

The authors have no relevant financial or non-financial interests to disclose.

\subsection*{Data Availability}

The cyclone tracks used in this study were obtained from the \textit{Atlantic Extratropical Cyclone Tracks Database}, available at \url{https://doi.org/10.17632/kwcvfr52hp.4}. All datasets, including the computed Lorenz Energy Cycle results for each cyclone, along with the scripts used to generate the figures and analyses presented in this study, are publicly available at the following GitHub repository: \url{https://github.com/daniloceano/energetic_patterns_cyclones_south_atlantic}. The source code of the \texttt{Cyclophaser} package, used for detecting cyclone life cycle phases, is available on PyPI at \url{https://pypi.org/project/cyclophaser/}. The \texttt{LorenzCycleToolkit} package, used for computing the Lorenz Energy Cycle components, is also available on PyPI at \url{https://pypi.org/project/LorenzCycleToolkit/}.

\subsection*{Author contributions}

Conceptualization: Danilo Couto de Souza, Pedro Leite da Silva Dias, Ricardo de Camargo, Carolina Barnez Gramcianinov; Methodology: Danilo Couto de Souza, Pedro Leite da Silva Dias, Ricardo de Camarg; Formal analysis and investigation: Danilo Couto de Souza, Pedro Leite da Silva Dias; Writing - original draft preparation: Danilo Couto de Souza; Writing - review and editing: Danilo Couto de Souza, Pedro Leite da Silva Dias, Ricardo de Camargo, Carolina Barnez Gramcianinov; Funding acquisition: Pedro Leite da Silva Dias, Ricardo de Camargo; Resources: Pedro Leite da Silva Dias, Ricardo de Camargo; Supervision: Pedro Leite da Silva Dias, Ricardo de Camargo.

\bibliography{sn-bibliography}% common bib file
%% if required, the content of .bbl file can be included here once bbl is generated
%%\input sn-article.bbl


\end{document}
